\def\fn#1{\em $\langle$#1$\rangle$}
\def\BL{\noindent $\bullet$ }
\def\carrot{$\wedge$}
% \documentclass [12pt] {article}     % Specifies the document class
  \documentclass [12pt] {ltxdoc}      % Specifies the document class

  \usepackage{graphics}
  \usepackage{amsmath}

                        \textwidth  16cm
                        \textheight 22cm
                        \hoffset=-2.0cm
                        \voffset=-2.0cm
\usepackage{titlesec}
\setcounter{secnumdepth}{4}
\titleformat{\paragraph}
{\normalfont\normalsize\bfseries}{\theparagraph}{1em}{}
\titlespacing*{\paragraph}
{0pt}{3.25ex plus 1ex minus .2ex}{1.5ex plus .2ex}

\begin{document}

\font\fcnfont=cmss10   % fn = function
\font\varfont=cmssi10  % vr = variable
\def\valfont{\tt}      % va = value of variable
\def\vshfont{\tt}      % vs = shell variable
\def\filfont{\rm}      % fl = file name
\def\txtfont{\tt}      % tx = text
\font\recfont=cmss12   % rc = G2F record
\def\cmdfont{\tt}      % cm = command
\def\prmfont{\tt}      % pa = parameter
\def\prifont{\bf}      % pr = primitive

\def\scn#1{Section~\ref  {#1}
               (page \pageref {#1})}
\def\scnn#1{Section~\ref  {#1},
               page \pageref {#1}}
\def\fcn#1{\hbox{\fcnfont #1}
          \def\_{-}
          (Section~\ref  {fcn-#1},
          page \pageref {fcn-#1})\def\_{\leavevmode
          \kern.06em \vbox{\hrule width0.3em}}}
\def\vr#1{{\hyphenpenalty=10000{\varfont #1}}}
\def\va#1{{\hyphenpenalty=10000{\valfont #1}}}
\def\fn#1{{\hyphenpenalty=10000{\fcnfont #1}}}
\def\fl#1{{\hyphenpenalty=10000{\filfont #1}}}
\def\tx#1{{\hyphenpenalty=10000{\txtfont #1}}}
\def\rc#1{{\hyphenpenalty=10000{\recfont #1}}}
\def\cm#1{{\hyphenpenalty=10000{\cmdfont #1}}}
\def\vs#1{{\hyphenpenalty=10000{\vshfont #1}}}
\def\pa#1{{\hyphenpenalty=10000{\prmfont #1}}}
\def\pr#1{{\hyphenpenalty=10000{\prifont #1}}}

% ****************
% ** Title Page **
% ****************

\centerline {\LARGE \bf RSIM MODELS AND ALGORITHMS}

\centerline {Adam N. Rosenberg}

\centerline {2022 November 15}

{\footnotesize
    \pagenumbering {roman}
    \tableofcontents
    \vfill \eject
    \pagenumbering {arabic}
}

\section {INTRODUCTION}
  \label {INTRODUCTION}

We can ask the question how retail unit sales depend on time.
We're pretty sure that higher prices lower the number sold,
at least most of the time,
but we want to quantify the relationship.
We could use a straight line or
we could a constant-elasticity model,
but, instead, we use a model where elasticity
is a constant times price,
\begin{equation}
E = \beta P.
\end{equation}
Without giving away too much of the plot,
I'll point out the solution of that differential equation is
\begin{equation}
Q = Q_0 e^{-\beta P} .
\end{equation}

So how do we find $Q_0$ and $\beta$?
Mostly we care about response to price $\beta$
rather than the quantity $Q_0$
as we're in the game of setting prices.
I sometimes write these relationships as
\begin{equation}
E = \frac{P}{\alpha}
\end{equation}
and
\begin{equation}
Q = Q_0 e^{P/\alpha}
\end{equation}
where $\alpha = 1/\beta$ is in the units of money
and has a physical representation of being what I call
{\em willingness to pay}.

All of that may be interesting,
but the form that is most amenable to the calculations
in this write-up is
\begin{equation}
Q = e^{q_0-\beta P}
\end{equation}
where $q_0 = \log Q_0$.

In this version of the story,
we're going to start with a history of $N$ time periods
$i=1,2,...,n$ with some price $P_1,P_2,...,P_N$ and
observations $X_1,X_2,...,X_N$.
Our objective is to get the truest, best, most-accurate forecast
and, along the way, to get the best $\beta$ response to price
for our model.

\section {SINGLE NORMAL DISTRIBUTION}
  \label        {NORMAL}

The simplest, easiest, most pedigogical method
is minimizing least squared error
(LSSE~=~least sum of squared error).
Let's assume we're fitting a linear equation
\begin{equation}
Q = a+bP
\end{equation}
instead of our more-complicated equation,
where $P$ is our independent variable,
$a$ and $b$ are coefficients we seek,
and we have sample data $P_1,P_2,...,P_N$ and $X_1,X_2,...,X_N$.

We measure the error in outcome
\begin{equation}
SSE = \sum_{i=1}^N (Q_i - X_i)^2 = \sum_{i=1}^N (a+bP_i-X_i)^2,
\end{equation}
take its derivative, and set that derivative equal to zero.
(There's some business of assuring
a positive-definite second derivative,
but we'll put that on a back burner.)
\begin{equation}
\frac{\partial SSE}{\partial a} =
\sum_{i=1}^N 2(a+bP_i-X_i) = 0
\end{equation}
\begin{equation}
\frac{\partial SSE}{\partial b} =
\sum_{i=1}^N 2(a+bP_i-X_i)P_i = 0
\end{equation}
We can separate and combine terms to get
\begin{equation}
    \begin{bmatrix}
                   N &
        \sum_{i=1}^N P_i \\
        \sum_{i=1}^N P_i &
        \sum_{i=1}^N P_i^2
    \end{bmatrix}
    \begin{bmatrix}
        a \\
        b
    \end{bmatrix}
     =
    \begin{bmatrix}
        \sum_{i=1}^N X_i \\
        \sum_{i=1}^N X_i P_i
    \end{bmatrix}
\end{equation}
and I'll leave solving that system of equations
as an exercise for the reader
as we're not using LSSE in our statistical solution of $\beta$.

A simpler application of LSSE is simply $Q=a$ without
any dependence on $P$.
\begin{equation}
SSE = \sum_{i=1}^N (Q_i - X_i)^2 = \sum_{i=1}^N (a-X_i)^2
\end{equation}
\begin{equation}
\frac{\partial SSE}{\partial a} =
\sum_{i=1}^N 2(a-X_i) = 0
\end{equation}
\begin{equation}
a=\sum_{i=1}^N \frac{X_i}{N}.
\end{equation}
It's intuitively satisfying that the constant best summarizing
a collection of samples is the mean value,
but a lot of assumptions go into that comfortable satisfaction.

The good news is
the LSSE technique is simple and requires little
advanced mathematics to solve.
The bad news is that LSSE imposes a scale on the value of the error.
If $(Q_1-X_1)^2$ goes down more than $(Q_2-X_2)^2$  goes up,
then we consider the change a win,
no matter how major or minor the actual change is.
If $Q_1$ and $X_1$ are around $0.1$
and $Q_2$ and $X_2$ are around $100.0$,
then I might not feel that way.
In any case,
just as economic elasticity removes the implicit valuation
of currency and quantity units,
the technique of maximum likelihood removes the implicit valuation
of scale.

Let's assume we have the same collection of observations
with the model assumption that the distributions of all
the normal distributions around $a$ have the same
standard deviation $\sigma$.
After a flurry of hand waving about conditional probabilities
we ask the likelihood of seeing $X_1,X_2,...,X_N$
as normally distributed samples of $N(a,\sigma^2)$
where $N(\mu,\sigma^2)$ is normal
with mean $\mu$ and variance $\sigma^2$.
Likelihood as a function of $a$ is
\begin{equation}
L = L(a) = \prod_{i=1}^N \frac{1}{ \sigma \sqrt{2 \pi}}
                     e^{-\frac{(a-X_i)^2}{2\sigma^2}}
\end{equation}
and we want to find $a$ for the minimum value of $L(a)$.
Solving a product like that is a mess,
so we take the natural logarithm of both sides.
We can get away with that for a minimization
because the natural logarithm is strictly and smoothly increasing.
\begin{equation}
\mathcal {L} = \mathcal {L}(a) =
    \log_e \frac{1}{ \sigma \sqrt{2 \pi}}
    \sum_{i=1}^N -\frac{(a-X_i)^2}{2\sigma^2}
\end{equation}
Aggregating the various constant terms we get a simple equation.
\begin{equation}
\mathcal {L} = \mathcal {L}(a) = K \sum_{i=1}^N (a-X_i)^2
\end{equation}
whose minimum-$\mathcal{L}$ solution is
\begin{equation}
a=\sum_{i=1}^N \frac{X_i}{N}.
\end{equation}

While this example may seem pedestrian and silly, or stupid,
I'll point out it also works, less elegantly but correctly,
if there are different weights and standard deviations.

\section {SINGLE POISSON DISTRIBUTION}
  \label        {POISSON}

Okay, we got through a single normal distribution
estimating the mean $a$ with constant variance $\sigma^2$.
Our next step is a Poisson distribution.
\begin{equation}
P(X \mid \lambda) = \frac {e^{-\lambda} \lambda^X}{X!}
\end{equation}
for integer $X \geq 0$
which is the distribution of all kinds of statistical stuff.
In particular, for retail,
it's the limiting distribution of a large number of shoppers $M$
each of whom has a small probability $p$ of buying one unit
with the limiting product $\lambda = Mp$.
For our illustrative example,
let's assume $\lambda$ is constant and we want to estimate it
from a sequence of $N$ observations $X_1,X_2,...,X_N$.

Let's walk through the process starting with the
likelihood function $L$,
taking the natural logarithm of both sides,
and setting the derivative equal to zero.
\begin{equation}
L(\lambda) = \prod_{i=1}^N
\frac {e^{-\lambda} \lambda^X_i}{X_i!}
\end{equation}
\begin{equation}
\mathcal{L}(\lambda) = \sum_{i=1}^N
-\lambda + X_i \log_e (\lambda) + \log_e {X_i!}
\end{equation}
\begin{equation}
\frac{\partial\mathcal{L}}{\partial\lambda} =
\sum_{i=1}^N -\lambda + \frac{X_i}{\lambda} = 0
\end{equation}
\begin{equation}
\lambda = \frac{\sum_{i=1}^N X_i}{N}
\end{equation}

I'll point out this is a robust solution.
I remember suggesting that a change of variable
would change the answer,
that this was just another form of min-error fitting.
My friend waited patiently as I substituted $\mu^2$ for $\lambda$
and re-did all the algebra only to get
\begin{equation}
\mu^2 = \frac{\sum_{i=1}^N X_i}{N}
\end{equation}
instead.

Up until here I've simply been exposing the reader to tools
for statistical analysis that go beyond scale in either variable
and that work on the Poisson distribution.
Now we're ready for the smallest piece of the big time.

\section {POISSON DISTRIBUTION WITH PRICE}
  \label                           {PRICE}

Suppose we have our original model from the introduction,
\begin{equation}
Q = e^{q_0 - \beta P}
\end{equation}
with a history of $N$ time periods
$i=1,2,...,n$ with some price $P_1,P_2,...,P_N$ and
observations $X_1,X_2,...,X_N$.

Two things make me cringe from any kind of metric-based fit
like least sum of squared error (LSSE).
First is a sense of scale,
when do big errors in big numbers win out
over small errors in small numbers?
Second, any scale looks wrong with the exponential
and taking the logarithm starts to get scary
when the logarithm of zero gets involved.
So, after our two easy rehearsals,
let's look at maximum likelihood.
\begin{equation}
L = L(q,\beta)
= \prod_{i=1}^N \frac{e^{-Q_i} Q_i^{X_i}}{X_i!}
= \prod_{i=1}^N \frac{e^{-e^{q_0 - \beta P_i}}
                     (e^{q_0 - \beta P_i})^{X_i}}{X_i!}
\end{equation}
\begin{equation}
\mathcal{L} = \mathcal{L}(q,\beta)
= \sum_{i=1}^N -Q_i + \log_e (Q_i) X_i - \log_e X_i!
= \sum_{i=1}^N -e^{q_0 - \beta P_i} + (q_0 - \beta P_i) X_i
                                           - \log_e X_i!
\end{equation}

I did this the long way using the far-right formula versions,
but I'm going to make life quite a bit simpler by using
the formulations with $Q$ and its derivatives.
\begin{equation}
\frac{\partial Q}{\partial q} = Q
\end{equation}
\begin{equation}
\frac{\partial Q}{\partial \beta} = -PQ
\end{equation}
\begin{equation}
\frac{\partial \log_e Q}{\partial q} = 1
\end{equation}
\begin{equation}
\frac{\partial \log_e Q}{\partial \beta} = -P
\end{equation}

So here we go:
\begin{equation}
\frac{\partial \mathcal{L}}{\partial q} =
  \sum_{i=1}^N -Q_i + X_i = 0
\end{equation}
\begin{equation}
\frac{\partial \mathcal{L}}{\partial q} =
  \sum_{i=1}^N P_i Q_i - P_i X_i = 0
\end{equation}
Isn't that sweet?
The first is that $q$ and $\beta$ are chosen
so total forecast unit sales equals total actual unit sales
and the second is that
so total forecast revenue equals total actual revenue.

Alas, finding those values isn't trivial.
Isaac~Newton and Joseph~Raphson were kind enough
to give us a way of finding zeroes of a differentiable function,
even one in more than one dimension.
If we're looking for a solution of $F(X)=0$
(not the same $X$ as before),
then we make an initial guess $X_0$ and recurse
\begin{equation}
X_{k+1} = X_k - F'(X_k)^{-1}(F(X_k)
\end{equation}
which, in one dimension, we can write
\begin{equation}
X_{k+1} = X_k - \frac{F(X_k)}{F'(X_k)}.
\end{equation}
We can think of this as ``sliding down the tangent line'' at $X_k$.

If we're solving $\mathcal{L}'(q,\beta)=0$,
then our iteration looks like
\begin{equation}
    \begin{bmatrix} q_{k+1} \\
                   \beta_{k+1} \end{bmatrix}
     =
    \big( \mathcal{L}''(q_k,\beta_k) \big)^{-1}
    \begin{bmatrix} q_{k} \\
                   \beta_{k} \end{bmatrix}
     =
    \begin{bmatrix}
        \frac {\partial^2 \mathcal{L}}{\partial q^2}(q_k, \beta_k) &
        \frac {\partial^2 \mathcal{L}}{\partial  q
                                       \partial \beta}(q_k, \beta_k)
          \\
        \frac {\partial^2 \mathcal{L}}{\partial \beta
                                       \partial  q}(q_k, \beta_k) &
        \frac {\partial^2 \mathcal{L}}{\partial \beta^2}(q_k, \beta_k)
    \end{bmatrix}^{-1}
    \begin{bmatrix} q_{k} \\
                   \beta_{k} \end{bmatrix}
\end{equation}
The $\mathcal{L}''$ second-derivative matrix
is called the Hessian ($H$).

Computing the Hessian goes like this.
\begin{equation}
H = \mathcal{L}''(q,\beta)
  = \begin{bmatrix}
        \frac {\partial^2 \mathcal{L}}{\partial q^2}(q_k, \beta_k) &
        \frac {\partial^2 \mathcal{L}}{\partial  q
                                       \partial \beta}(q_k, \beta_k)
          \\
        \frac {\partial^2 \mathcal{L}}{\partial \beta
                                       \partial  q}(q_k, \beta_k) &
        \frac {\partial^2 \mathcal{L}}{\partial \beta^2}(q_k, \beta_k)
    \end{bmatrix}
  = \begin{bmatrix}
        - \sum_{i=1}^N       Q_i &
          \sum_{i=1}^N P_i   Q_i \\
          \sum_{i=1}^N P_i   Q_i &
        - \sum_{i=1}^N P_i^2 Q_i
    \end{bmatrix}
\end{equation}

There's one other cool thing here.
One of the cool, lucky things about this problem is
in solving a maximum likelihood problem of a Poisson distribution
of an exponential of a linear function in one or more variables
is that the inverse of the Hessian ($H^{-1}$),
the same inverse we compute for each Newton-Raphson step,
is precisely the covariance matrix we need
to calculate standard deviations of our coefficients.

This is the state of the retail price-modeling art
as of 2011 when we incorporated Clear~Demand.
The $P_i$ and $X_i$ came from input sales data
when $X_i > 0$ and are estimates from a clever
out-of-stock algorithm for time periods
where a product and store have no sales at all.

\section {TRUNCATED POISSON DISTRIBUTION}
  \label {TRUNCATED}

Two things occurred to me
when considering the issue of out-of-stock time periods
for a product and store.
First, we don't really know which of the zero-sales time periods
are out of stock (where the product was not available for sale)
or in stock (where the shoppers didn't buy any).
If the product is in stock for that time period,
then we don't know the actually price shoppers did not pay that day.
Second, the model depends significantly on the results
of any out-of-stock projects we might do based on gaps in unit sales.
So wouldn't it be wonderful if we model demand
independent of uncertainty about zero-sales time periods?

Consider a product with sales averaging twenty-five units per day
at a particular store.
The Poisson distribution with mean $\lambda=25.0$
has variance $\sigma^2=25.0$, standard deviation $\sigma=5.0$,
and probabality 0.000000002 of being zero.
(That's once in 500~million samples, $e^{-25}$.)
If I make a bald, sweeping assumption that all the zero-sales samples
are out of stock, or otherwise unavailable,
then I'm hardly ever going to be wrong.
Even if I relax my assumption and figure
a standard deviation of $\sigma=10.0$ for a mean of $\mu=25.0$
I'll only be wrong once in 160~samples, not terribly far off.

When the average sales gets down to two per day $\lambda=2.0$
things get squishy.
Blocks of time with no sales might be out of stock,
but they also may be legitimate zero-sales days.
The vector of probabilities for $\lambda=2.0$ starts with
\begin{align}
        Prob (X=0) &= 0.1353 \nonumber \\
        Prob (X=1) &= 0.2706 \nonumber \\
        Prob (X=2) &= 0.2706 \nonumber \\
        Prob (X=3) &= 0.1804           \\
        Prob (X=4) &= 0.0902 \nonumber \\
        Prob (X=5) &= 0.0361 \nonumber \\
        Prob (X=6) &= 0.0120 \nonumber
\end{align}
so about 14\% of the days there are zero sales,
27\% of the days there is one unit sold,
another 27\% of the days two units are sold,
18\% of the days sell three units,
{\em et cetera}.
Our problem is that some of the zeros are ``real''
in that they are legitimate nothing-sold days
and others are out of stock, unavailable, off the shelf,
and we have no reliable way of telling which is which.

Back in Khimetrics days we estimated out of stock
by using the average sales $r$ and estimating probabilities
of sequences of $n$ zeros as $e^{-nr}$.
When $e^{-nr}$ fell below some threshold,
usually 0.02,
then we figured the outage was out of stock
rather than available but not sold.
A lot of our estimate process of both $q$ and $\beta$
are highly swayed by changes in out-of-stock estimation.

Here's my idea (patent pending, I hope).
Look at the {\em truncated} distribution,
the Poisson conditioned on at least one unit sold.
We divide by $1-e^{-\lambda}$ to get conditional probabilities
for the same Poisson distribution.
\begin{align}
        Prob (X=1 \mid X>0)&= 0.3130 \nonumber \\
        Prob (X=2 \mid X>0)&= 0.3130 \nonumber \\
        Prob (X=3 \mid X>0)&= 0.2086 \nonumber \\
        Prob (X=4 \mid X>0)&= 0.1043           \\
        Prob (X=5 \mid X>0)&= 0.0417 \nonumber \\
        Prob (X=6 \mid X>0)&= 0.0139 \nonumber
\end{align}
In the case of $\lambda=2.0$ there's still some ``meat on the bone''
even without the zero values.
Consider the case of $\lambda=0.2$ and the picture is different.
\begin{align}
        Prob (X=0) &= 0.8187  \nonumber \\
        Prob (X=1) &= 0.1637  \nonumber \\
        Prob (X=2) &= 0.0164            \\
        Prob (X=3) &= 0.0011  \nonumber \\
        Prob (X=4) &= 0.00005 \nonumber
\end{align}
There are still some samples with $X>1$, but not so many.
Our truncation to eliminite $X=0$ points yields this vector
with only one-fifth as many points ($Prob (X>0)=0.1813$ above).
\begin{align}
        Prob (X=0) &= 0.9033 \nonumber \\
        Prob (X=1) &= 0.0903 \nonumber \\
        Prob (X=2) &= 0.0060           \\
        Prob (X=3) &= 0.0030 \nonumber \\
        Prob (X=4) &= 0.0001 \nonumber
\end{align}
We're basing our entire estimate of $\lambda$
on 10\% of the time periods $X>0$
which are 2\% of the total.

By the time we get down to $\lambda=0.02$,
a really slow selling product,
98\% of the available time periods are zero sales
and removing them from consideration
is a major cut of useful information.
Now 99\% of our remaining time periods have $X=1$
and we're basing our entire estimate of $\lambda$
on one percent of the non-zero time periods.
Gosh, if only we could have the zeroes back
with some kind of confidence about their availability status!

Well, up until now, I believe we aren't going to get that confidence.
We really don't know the statistical status
of products and stores with $\lambda < 0.5$
and our estimates of them really are soft and squishy.
I believe trying to guess out-of-stock {\em versus} none-sold days
creates an appearance of knowing what we really don't know.

So let's explore the same maximum-log-likelihood model
with the extra twist that we only care about
days with {\em positive} sales, not the zeros.
That means we have a {\em truncated} Poisson distribution
\begin{equation}
P(X \mid \lambda) = \frac {e^{-\lambda} \lambda^X}{X!
    (1-e^{-\lambda})}
\end{equation}
for integer $X > 0$ (not~zero).
That gives us the following similar footsteps.
\begin{equation}
L = L(q,\beta)
= \prod_{i=1}^N \frac{e^{-Q_i} Q_i^{X_i}}{X_i! (1-e^{-Q_i})}
= \prod_{i=1}^N \frac{e^{-e^{q_0 - \beta P_i}}
         (e^{q_0 - \beta P_i})^{X_i}}{X_i! (1-e^{-e^{q_0 - \beta
         (P_i}})}
\end{equation}
\begin{align}
\mathcal{L} = \mathcal{L}(q,\beta)
&= \sum_{i=1}^N -Q_i + \log_e (Q_i) X_i - \log_e X_i!
                                       - \log_e (1-e^{-Q_i}) \\
&= \sum_{i=1}^N -e^{q_0 - \beta P_i} + (q_0 - \beta P_i)X_i -\log_e
&= X_i!
                                 - \log_e (1-e^{-e^{q_0 - \beta P_i}})
\end{align}

So here we go again:
\begin{equation}
\label{UnitsEquation}
\frac{\partial \mathcal{L}}{\partial q} =
  \sum_{i=1}^N -Q_i + X_i - \frac{1}{(1-e^{-Q})}(e^{-Q})Q = 0
\end{equation}
\begin{equation}
\label{RevenueEquation}
\frac{\partial \mathcal{L}}{\partial \beta} =
  \sum_{i=1}^N P_i Q_i - P_i X_i + \frac{1}{(1-e^{-Q})}(e^{-Q})PQ = 0
\end{equation}

\begin{align}
H &= \begin{bmatrix}
        \frac {\partial^2 \mathcal{L}}{\partial q^2}(q_k, \beta_k) &
        \frac {\partial^2 \mathcal{L}}{\partial  q
                                       \partial \beta}(q_k, \beta_k)
          \\
        \frac {\partial^2 \mathcal{L}}{\partial \beta
                                       \partial  q}(q_k, \beta_k) &
        \frac {\partial^2 \mathcal{L}}{\partial \beta^2}(q_k, \beta_k)
    \end{bmatrix} \\
  &= \begin{bmatrix}
          \sum_{i=1}^N -       Q_i
                             +       \frac   {Q_i
          e^{-Q_i}}{1-e^{-Q_i}}
                              \left( \frac   {Q_i
          e^{-Q_i}}{1-e^{-Q_i}}
                                     - Q +1 \right) &
          \sum_{i=1}^N P_i \left( Q_i
                             -       \frac   {Q_i
          e^{-Q_i}}{1-e^{-Q_i}}
                              \left( \frac   {Q_i
          e^{-Q_i}}{1-e^{-Q_i}}
                                     - Q +1 \right) \right) \\
          \sum_{i=1}^N P_i     Q_i
                             -       \frac   {Q_i
          e^{-Q_i}}{1-e^{-Q_i}}
                              \left( \frac   {Q_i
          e^{-Q_i}}{1-e^{-Q_i}}
                                     - Q +1 \right) &
          \sum_{i=1}^N P_i^2 \left( - Q_i
                             +      \frac    {Q_i
          e^{-Q_i}}{1-e^{-Q_i}}
                             \left( \frac    {Q_i
          e^{-Q_i}}{1-e^{-Q_i}}
                                     - Q +1 \right) \right)
    \end{bmatrix}
\end{align}
For those who like bigger print, I can write it this way.
\begin{align}
        \frac {\partial^2 \mathcal{L}}{\partial q^2}(q, \beta)
    &=    \sum_{i=1}^N -       Q_i
                             +       \frac   {Q_i
          e^{-Q_i}}{1-e^{-Q_i}}
                              \left( \frac   {Q_i
          e^{-Q_i}}{1-e^{-Q_i}}
                                     - Q +1 \right) \\
        \frac {\partial^2 \mathcal{L}}{\partial  q
                                       \partial \beta}(q, \beta)
    =   \frac {\partial^2 \mathcal{L}}{\partial \beta
                                       \partial  q}(q, \beta)
   &=     \sum_{i=1}^N P_i \left( Q_i
                             -       \frac   {Q_i
          e^{-Q_i}}{1-e^{-Q_i}}
                              \left( \frac   {Q_i
          e^{-Q_i}}{1-e^{-Q_i}}
                                     - Q +1 \right) \right) \\
        \frac {\partial^2 \mathcal{L}}{\partial \beta^2}(q, \beta)
   &=     \sum_{i=1}^N P_i^2 \left( - Q_i
                             +      \frac    {Q_i
          e^{-Q_i}}{1-e^{-Q_i}}
                             \left( \frac    {Q_i
          e^{-Q_i}}{1-e^{-Q_i}}
                                     - Q +1 \right) \right)
\end{align}

\section {COMPRESSING THE Q's}
  \label {COMPRESSING}

Consider a collection of products in a price family
and stores in a market.
In our abstraction-based world,
we call the latter {\em locations} in an {\em area}.
This Cartesian product in product-store space
is what we call a {\em model group}
and the entire model group gets a single $\beta$ value
while each product and location have their own $q$ value.
If a price family has ten products and
an area has twenty locations (market with twenty stores),
then we have 200~$q$ values and one $\beta$ in
a 201-by-201 system of equations.
As inverting a matrix is an effort roughly proportional
to the cube of the matrix size,
compressing a matrix by a factor of ten
reduces the inversion effort by a factor of roughly one thousand.

Consider a linear process where $X=a+bP+\epsilon$ where
$a$ and $b$ are constants, $P$ is the independent variable,
and $\epsilon$ is an unknowable error term normally distributed.
Given a set of sample data $P_1,P_2,...,P_n$ and  $X_1,X_2,...,X_N$
we can use our LSSE approach to find $a$ and $b$.
We know that already.

Now consider a slightly more complicated linear process
where $X=a_R+bR+\epsilon$ for $R_1,R_2,...,R_N$
and $X=a_S+bS+\epsilon$ for $S_{N+1},S_{N+2},...,S_{N+M}$
with the same $b$ and distinct constants $a_R$ and $a_S$.
We could solve a single system with three variables
$X=a_r r + a_s s +bP$ where $P$ gets the original $R$ and $S$ values
and the new $r$ variable is one for $R_1,R_2,...,R_N$ and zero
    otherwise
and the new $s$ variable is similarly
one for $S_{N+1},S_{N+2},...,S_{N+M}$ and zero otherwise.
If we have oodles of sets with distinct and different
$a$ constant coefficients all with the same slope $b$,
then we get a situation like having
lots of products in a price family and lots of locations in an area
with a Cartesian product of $q_{ProductLocation}$ values and
one solitary $\beta$ response to price.

Two things come to mind.
First, the matrices are getting larger
with larger price families and areas,
more computation for the modeling computer.
Second, the fraction of the effort that's useful for $\beta$
is shrinking.
Since we know the conditions of maximum likelihood include
\begin{equation}
\sum_{i=1}^N X_i =
\sum_{i=1}^N Q_i =
\sum_{i=1}^N e^{q-\beta P_i} =
e^q \sum_{i=1}^N e^{-\beta P_i},
\end{equation}
it isn't hard to find $q$ for each product and location
once we have $\beta$.
So spending 99\% of our computational energy
to find maybe one hundred $q$ values doesn't sound ideal.

Let's go back to our two-variable fit.
We could do the most aggressive thing and fit $Q=a+bP$
while we completely ignore the utter bimodelity of the
intercept variables $a_r$ and $a_s$.
Would that work?
Well, so long as the $P$ values for both $R$ and $S$ variables
come from the same distribution, then the slope we see and fit
{\em should} be the slope $b$.
(I'm too lazy to invoke a graphics package.)
\begin{minipage}{\linewidth}
\begin{verbatim}
                |                        S      S
                |                S    S    S
             X  |          S        S        S
                |     S        S                ___
                |        S            _____----
                |   S        _____----
                |   _____----                   R
                |               R         R      R
                |        R   R        R R
                |    R  R         R
                |
                +-------------------------------- R S
\end{verbatim}
\end{minipage}
That's asking a lot as even the slightest difference in sampling
range between $R$ and $S$ values gives us a false slope $b$.
\begin{minipage}{\linewidth}
\begin{verbatim}
                |                        S
                |                S    S    S
             X  |          S        S        S
                |   __S_       S
                |       -S--____
                |   S           ----____
                |                       ----____R
                |               R         R     -R-
                |            R        R R
                |                 R
                |
                +-------------------------------- R S
\end{verbatim}
\end{minipage}

If we have prior knowledge of the {\em difference} $D$
between $R$ and $S$ in this example,
then we combine the two and get our original, same slope $b$
solving for only two variables instead of three.
\begin{minipage}{\linewidth}
\begin{verbatim}
                |                        S      R ___
                |               RS    S __RS_----R
             X  |          S R ____S-R-R    S
                |     S____----S  R
                | ----   S
                |   S
                |
                |
                |
                |
                |
                +-------------------------------- R+D S
\end{verbatim}
\end{minipage}

We combine stores this way in our model
using overall difference in store levels
estimated from total sales.
Figuring a store whose total sales are twice another store
should forecast twice the amount,
we adjust the forecasts and combine $q$ values
for locations in an area.
We don't do this among products in a price family.
This compresses the problem without compromising
the statistical validity of the response to price $\beta$.
(What never?  Well hardly ever!)
Once we get $\beta$ from this location-combined model,
we go back to figure out the $q$ values by product-location.
\begin{equation}
q = \frac {\sum_{i=1}^N X_i}
          {\sum_{i=1}^N e^{-\beta P_i}}
\end{equation}

Consider the case where we're inferring $q$ from $\beta$
with no-zeros input-data as in the section above.
We don't know how many zero-sales days there are,
but a good guess is $1/(1-e^{-X})$ for a day with $X$ units.
\begin{equation}
q = \frac {\sum_{i=1}^N \frac{X_i}{1-e^{-X_i}}}
          {\sum_{i=1}^N e^{-\beta P_i}}
\end{equation}

                \subsection {Solving For Q}
                     \label {Solving}

I thought the solution of $q$ from $\beta$ and sales data
was the easy part.
The maximum likelihood equilibrum condition in
Equations~\ref {UnitsEquation} tells us
\begin{equation}
\sum_{i=1}^N Q_i = \sum_{i=1}^N X_i
\end{equation}
where $Q_i=e^{q_0-\beta P_i}$
and $X_i$ are the actual sales observations.
Solving (mostly) for $q_0$ gives us
\begin{equation}
e^{q_0} \sum_{i=1}^N e^{-\beta P_i} = \sum_{i=1}^N X_i
\end{equation}
and then
\begin{equation}
e^{q_0} = \frac {\sum_{i=1}^N X_i}{ \sum_{i=1}^N e^{-\beta P_i}}
\end{equation}
and finally
\begin{equation}
Q_0 = \frac {\sum_{i=1}^N X_i}{ \sum_{i=1}^N e^{-\beta P_i}}
\end{equation}
where $Q_0=e^{q_0}$.

Let's explore that in more detail.
Let $f_i$ be the $\beta$ forecast term $e^{-\beta P}$,
the part we're treating
as fixed in our search for $q_0$ and $Q_0$.
It might even be $Se^{-\beta P}$ where $S$ is seasonality,
also fixed in our search for $q_0$ and $Q_0$.

Let's do the maximum likelihood with $X_i$
being a Poisson distribution of some quantity $Q_i$
where $Q_i = Q_0 f_i$.
As before $L$ is the likelihood function
and $\mathcal{L}$ is the log of the likelihood function.
\begin{equation}
L = L(Q) = \prod_{i=1}^n \frac {e^{-Q_i} Q_i^{X_i}}{X_i !}
         = \prod_{i=1}^n \frac {e^{Q_0 f_i} (Q_0 f_i)^{X_i}}{X_i !}
\end{equation}
\begin{equation}
\mathcal{L} = \mathcal{L}(Q) =
     \sum_{i=1}^n -Q_i + \log_e (Q_i) X_i - \log_e X_i !
   = \sum_{i=1}^n -Q_0 f_i + \log_e (Q_0 f_i) X_i - \log_e X_i !
\end{equation}
\begin{equation}
\mathcal{L} = \mathcal{L}(Q)
   = \sum_{i=1}^N -Q_0 f_i
   + \log_e (Q_0) \sum_{i=1}^n X_i
   + \sum_{i=1}^N \log_e (f_i) X_i
   - \sum_{i=1}^N \log_e X_i !
\end{equation}
\begin{equation}
\mathcal{L} = \mathcal{L}(Q)
   = \sum_{i=1}^N -Q_0 f_i
   + \log_e (Q_0) \sum_{i=1}^N X_i
   + \sum_{i=1}^N \log_e (f_i) X_i
   - \sum_{i=1}^N \log_e X_i !
\end{equation}
\begin{equation}
\frac{\partial \mathcal{L}}{\partial Q_0} =
  - \sum_{i=1}^N f_i + \frac {\sum_{i=1}^N X_i}{Q_0} = 0
\end{equation}
\begin{equation}
   \sum_{i=1}^N f_i = \frac {\sum_{i=1}^N X_i}{Q_0}
\end{equation}
\begin{equation}
   Q_0 = \frac {\sum_{i=1}^N X_i}{\sum_{i=1}^N f_i}
\end{equation}
So the obvious answer is $Q_0$ is total units sold
divided by
total units forecast with price and maybe seasonality.

%====

Apparently this isn't as obvious as I thought.
One of my colleagues suggested doing
a least-sum-of-squared-error (LSSE) calculation for Q.
\begin{equation}
\mathrm{Minimize} ~ y = \sum_{i=1}^N (X_i - Q_i)^2
\end{equation}
% \begin{equation}
% \mathrm{Minimize} ~ y = \sum_{i=1}^N (X_i - Q_0 e^{-\beta P_i})^2
% \end{equation}
% Let $f$ be the $\beta$ forecast term.
% \begin{equation}
% f_i = e^{-\beta P_i}
% \end{equation}
\begin{equation}
\mathrm{Minimize} ~ y = \sum_{i=1}^N (X_i - Q_0 f_i)^2
\end{equation}
Take the derivative with respect to $Q_0$ and set to zero.
\begin{equation}
\frac {\partial y}{\partial Q_0}
                 = \sum_{i=1}^N 2 f_i (X_i - Q_0 f_i) = 0
\end{equation}
which gives us
\begin{equation}
\sum_{i=1}^N f_i X_i = \sum_{i=1}^N Q_0 f_i^2
\end{equation}
\begin{equation}
Q_0 = \frac {\sum_{i=1}^N f_i X_i }{ \sum_{i=1}^N f_i^2 }
\end{equation}

If we make these weighted averages with weights $w_i$,
then we get these for maximum likelihood and LSSE.
They're not the same.
The second formula (LSSE) gives greater weight
to lower-priced time periods.
\begin{equation}
\label{maxl01}
Q_0 = \frac {\sum_{i=1}^N X_i w_i}{ \sum_{i=1}^N f_i w_i}
~ ~ ~ ~ ~ ~ \mathrm {maximum~likelihood}
\end{equation}
\begin{equation}
\label{lsse01}
Q_0 = \frac {\sum_{i=1}^N X_i f_i w_i }{ \sum_{i=1}^N f_i^2 w_i}
~ ~ ~ ~ \mathrm {least~sum~of~squared~error}
\end{equation}
Before we get into too big a snit over this difference,
recall that the $\beta$ coefficient has been carefully chosen
to minimize the variation of $X_i / f_i$,
especially as a function of price as $f_i = e^{-\beta P_i}$.
So I don't think Equations~\ref {maxl01} and~\ref {lsse01}
should yield terribly different values of $Q_0$.

I'm not sure what to do about the questionable zero-sales events
in the least-sum-of-squares case,
but let's follow through on the maximum likelihood for $Q_0$.
\begin{equation}
L = L(Q_0) = \prod_{i=1}^N
    \frac{e^{-Q_i}Q_i^{X_i}} {X_i! (1-e^{-Q_i})}
\end{equation}
\begin{equation}
L = L(Q_0) = \prod_{i=1}^N
    \frac{e^{-Q_0 f_i} (Q_0 f_i)^{X_i}} {X_i! (1-e^{-Q_0 f_i})}
\end{equation}
\begin{equation}
\mathcal{L}=\mathcal{L}(Q_0) =
    \sum_{i=1}^N -Q_i + (\log_e Q_i) X_i
                  - \log_e X_i! - \log_e (1 - e^{Q_i})
\end{equation}
\begin{equation}
\mathcal{L}=\mathcal{L}(Q_0) =
    \sum_{i=1}^N -Q_0 f_i + (\log_e Q_0 f_i) X_i
                  - \log_e X_i! - \log_e (1 - e^{Q_0 f_i})
\end{equation}
\begin{equation}
\frac{\partial \mathcal{L}}{\partial Q_0} =
    - \sum_{i=1}^N f_i + \frac {1}{Q_0} \sum_{i=1}^N X_i
    - \sum_{i=1}^N \frac {1}{1-e^{-Q_0 f_i}} (e^{-Q_0 f_i}) f_i = 0
\end{equation}
\begin{equation}
\frac{\partial \mathcal{L}}{\partial Q_0} =
    - \sum_{i=1}^N f_i + \frac {1}{Q_0} \sum_{i=1}^N X_i
    - \sum_{i=1}^N (1-\frac{1}{1-e^{Q_0 f_i}}) f_i = 0
\end{equation}
\begin{equation}
\frac{\partial \mathcal{L}}{\partial Q_0} =
    - \sum_{i=1}^N f_i + \frac {1}{Q_0} \sum_{i=1}^N X_i
    - \sum_{i=1}^N f_i
    + \sum_{i=1}^N \frac{f_i}{1-e^{-Q_0 f_i}} = 0
\end{equation}
\begin{equation}
\label{iter01}
Q_0 = \frac {\sum_{i=1}^N X_i}
            {\sum_{i=1}^N f_i (2-\frac{1}{1-e^{-Q_0 f_i}})}
\end{equation}

We're not out of the woods completely
as there is that little $Q_0$ in the denominator of the denominator,
but that should iterate nicely to converge to a nice value for $Q_0$.
(Actually, I think the above equation with $1-e^{-X_i}$ on top
may be wrong, but it shouldn't be too far off and it's
easy to compute.)
We could start the process with
\begin{equation}
Q_0 = \frac {\sum_{i=1}^N \frac {X_i}{1-e^{-X_i}}}
            {\sum_{i=1}^N f_i}
\end{equation}
and then iterate by computing the right side
of Equation~\ref {iter01} until we get
a stable value of $Q_0$.
I don't think it will take a lot of iterations to settle down,
but I could be wrong.

% This still leaves the issue of
% out of stock and truncated distributions
% to be figured out.

% \centerline {==== STILL WRITING THIS PART ====}

\section {SEASONALITY}
  \label {SEASONALITY}

Suppose I know time-dependent effect on sales
other than price (or customer type or sales type).
I can use the same basic formulae fitting $Q-X$ and $PQ-PX$
knowing that the $Q_i$ have the extra time-dependent term,
which we call {\em seasonality}.
Once we have knowledge of seasonality $S$
we can set $Q$ to include it.
\begin{align}
Q &= SQ_0 e^{-\beta P} \\
Q_i &= S_i Q_0 e^{-\beta P_i}
\end{align}
Notice that $S$ can change from sample to sample $S_i$
even though it is not being searched by the statistical process
that we're using to hunt for $\beta$ and then for $q$.

First, let's separate our model universe into
{\em season groups} where each season group
is a collection of model groups whose seasonal behavior
is expected to be the same.
(Most of our actual cases are small-enough parts of the enterprise
that seasonality is consistent within them,
so we don't put patio furniture and snow shovels
in the same case.
Often enough to be interesting
we have distinct season groups within a single test case.)

Exponential of a linear polynomial is a good structure
for solving a maximum-log-likelihood problem.
\begin{equation}
Q = e^{A_0 + A_1 x_1 + A_2 x_2 + A_3 x_3 + A_4 x_4 + ...}
\end{equation}
In our case it looks like constant $q$ and
price response $\beta$ plus seasonality terms
$S_1, S_2, S_3, S_4, ...$,
$P$ is price, and $x_k$ is the seasonality variable value.
\begin{equation}
Q = e^{q - \beta P + S_1 x_1 + S_2 x_2 + S_3 x_3 + S_4 x_4 + ...}
\end{equation}
In our time series that looks like this
where $i=1...n$ is the time series.
\begin{equation}
Q_i = e^{q -\beta P_i + S_1 x_{i1} + S_2 x_{i2} +
                        S_3 x_{i3} + S_4 x_{i4} + ...}
\end{equation}
The maximum-log-likelihood solution described above
works the same with the extra variables,
especially if the terms are all linear as they are here.

The season variables are a time series for each variable
forming an array by time and season variable.
The first six Fourier terms for an annual-periodic function are
$ \mathrm {sin} ( \frac{date}{oneyear} \times 2\pi ) $,
$ \mathrm {cos} ( \frac{date}{oneyear} \times 2\pi ) $,
$ \mathrm {sin} ( \frac{date \times 2}{oneyear} \times 2\pi ) $,
$ \mathrm {cos} ( \frac{date \times 2}{oneyear} \times 2\pi ) $,
$ \mathrm {sin} ( \frac{date \times 3}{oneyear} \times 2\pi ) $, and
$ \mathrm {cos} ( \frac{date \times 3}{oneyear} \times 2\pi ) $.
We also have a trend term $(today-date)$.
All these terms are in the exponent
so there is an asymmetry in that $e_x-1 > 1-e^{-x}$
so this model's seasonal increases come out bigger
than its seasonal decreases.

We also have specific known holidays
with their own structure.
Most holidays are the same date every year,
probably with the same ramp-up before and ramp-down after,
but some religious holidays wander over the secular calendar
from year to year.
I'm thinking of the Christian Easter holiday
and almost all the Jewish holidays.
Also there may be holiday-type seasonal behavior
for Homecoming or the Senior Prom that have
the same behavior year to year,
then those should also be included,
especially if they don't always come at the same date.

Finally we have day of week, at least for daily models.
The lift (or drop) from Wednesday to Friday, for example,
is usually consistent and should be modeled as part of seasonality.
In theory, at least in this approach,
all these effects are additive in the exponent
which means they're multiplicative in actual sales,
so if September is average plus 20\%,
Labor~Day is average plus 15%,
and Friday is Wednesday plus 10\%
then the total lift should be 20\%+15\%+10\%,
$1.20 \times 1.15 \times 1.10$ which is $1.518$.

Fine, how do we calculate and estimate $S_i$?
Ultimately, the whole model universe,
or at least a whole season group,
is one or two dozen season-variable coefficients,
maybe a few hundred $\beta$ coefficients,
and a several thousand $q$ coefficients.
We compressed the $q$ coefficients
in Section~\ref {COMPRESSING},
but we still have lots of $\beta$ coefficients.
Once I have season-variable coefficients,
the $\beta$ calculations for each model group
are much simpler and more stable
than they would be all together.

So I figure I'll solve for seasonality-plus-$\beta$
in each model group and put them all together
to get season-variable coeffients.
A weighted average based on confidence and revenue
seems like a good way to get seasonality.
As these individual-model-group maximum-log-likelihood regressions
can get squirrely,
we'll put bounds on the season-variable coefficients.

Seasonality is probably the weakest link
in the maximum-log-likelihood demand model in this paper.
We have a collection of sort-of-right coefficients
being badly estimated repeatedly
in the hope that averaging them will get a good answer.
Since those individual-model-group estimates
can be ``all over the place,''
there's not much comfort that their average
represents true seasonality coefficients.
Furthermore, it's not clear that
the scope of seasonality estimation
should be the same as the other modeling.

% ================

\section {LOCAL OPTIMIZATION}
  \label {LOCAL OPTIMIZATION}

    (stuff about optimization and opportunity curve goes here)

Here are the basic equations for revenue ($R$)
and profit ($\Pi$, capital Greek letter $\pi$)
based on price ($P$), cost ($C$), and quantity ($Q$, unit sales).
\begin{align}
\label{eqn:opt01}
R   &= PQ \\
\Pi &= (P-C)Q
\end{align}
Using our model we have
\begin{align}
\label{eqn:opt02}
            Q &= Q_0 e^{-\beta P} \\
\frac{dQ}{dP} &= -\frac{1}{\beta} Q_0 e^{-\beta P} = -\beta Q.
\end{align}
We're trying to maximize value ($V$)
\begin{align}
\label{eqn:opt03}
V &= \lambda \Pi + (1-\lambda) R \\
V &= \lambda (P-C)Q + (1-\lambda) PQ \\
V &= (P-\lambda C)Q
\end{align}
Setting the derivative to zero we use the product rule to get
\begin{align}
\label{eqn:opt04}
\frac{dV}{dP} &= -(P-\lambda C)\beta Q + Q = 0 \\
          P^* &= \frac{1}{\beta} + \lambda C
\end{align}

                \subsection {Tax Included in Price}
                     \label {Tax}

    Suppose we have a tax that is a fraction ($\tau$)
of the price ($P$) that does not count as profit.
\begin{align}
\label{eqn:opt11}
  R &= PQ \\
\Pi &= (P-\tau P-C)Q \\
  V &= \lambda \Pi + (1-\lambda) R \\
V &= \lambda (P-\tau P-C)Q+(1-\lambda)PQ=(P-\lambda\tau P-\lambda C)Q \\
\frac{dV}{dP} &= -\beta(P-\lambda\tau P-\lambda C)Q+(1-\lambda\tau)Q=0
\end{align}
This last happens when
\begin{align}
\label{eqn:opt12}
(-\beta (P-\lambda\tau P-\lambda C)Q+(1-\lambda\tau)Q &= 0 \\
P(1-\lambda\tau)-\lambda C &= \frac{1}{\beta}(1-\lambda\tau) \\
P-\frac{\lambda C}{1-\lambda\tau} &= \frac{1}{\beta} \\
P &= \frac{1}{\beta}+\frac{\lambda}{1-\lambda\tau}C
\end{align}

Go along with me here.
Let's set a variable $\mu$
to $\lambda / (1-\lambda\tau)$.
\begin{align}
\label{eqn:opt13}
\mu &= \frac{\lambda}{1-\lambda \tau} \\
\mu-\lambda\tau\mu &= \lambda \\
\mu &= \lambda+\lambda\tau\mu \\
\label{eqn:opt14}
\lambda &= \frac{\mu}{1+\mu\tau}
\end{align}

Here's the cool part.
If we run our RSIM optimization as usual
and get a profit-revenue-weight $\lambda$
as our ``optimal'' value for $\lambda$,
and if there is a tax component $\tau P$ of price $P$,
then we can call that $\lambda$ value value $\mu$
and use Equation~\ref {eqn:opt14} to find
the profit-revenue weight $\lambda$ with the tax included.

Think of it this way.
Let's start with our opportunity curve with
a profit-revenue frontier.
Each point on that curve is associated with some
revenue-profit weight $\lambda$.

\begin{minipage}{\linewidth}
\begin{verbatim}
                |
                |
        profit  |
                |   oooooo
                |         ooooo
                |              oooo
                |                  ooo
                |                     oo
                |                       o
                |                       o
                +-------------------------------- revenue
\end{verbatim}
\end{minipage}

Now let's figure in the tax as a linear function
$\tau$ times revenue so the graph is stretched downward
as revenue increases.

\begin{minipage}{\linewidth}
\begin{verbatim}
                |
                |
        profit  |
                |   oo    
                |     oooo
                |         o     
                |          oooo       
                |              o       
                |               ooo
                |                  oo
                +-----               o
                      -----           o
                           -----       o
                                -----  o
                                     -----
                                          -----
                                               -- revenue
\end{verbatim}
\end{minipage}

Now let's put the horizontal-abscissa revenue axis
back where it belongs and check out
the new, displaced-downward curve.

\begin{minipage}{\linewidth}
\begin{verbatim}
                |
                |
        profit  |
                |   oo    
                |     oooo
                |         o     
                |          oooo       
                |              o       
                |               ooo
                |                  oo
                +--------------------o----------- revenue
                                      o
                                       o
                                       o
\end{verbatim}
\end{minipage}

Every point on the old profit-revenue frontier
is a point on the new profit-revenue frontier,
but with a different profit-revenue weight $\lambda$.
For a tax-included RSIM run
it looks like we can run RSIM the usual way,
call the resulting no-tax profit-revenue weights
$\mu$ instead of $\lambda$,
and use Equation~\ref {eqn:opt14}
to convert to the new, tax-included
profit-revenue weight $\lambda$.

% XXXXXXXXXXXXXXXXXXXXXXXXXXXXXXXXXXXXXXXXXXXXXXXXXXXXXXXXXXXXXXXX

\section {OVERDISPERSION}
  \label {OVERDISPERSION}

At the moment,
deep in our hearts,
we believe there are two major contributors
to variance $\sigma^2$ being larger than the
Poisson-indicated mean $\lambda$.
We call this larger-than-Poisson variance {\em overdispersion}.
Let's spend a moment going down to a more-primitive model of retail,
traffic, share, and count.
{\em Traffic} is the number of shoppers in a store or on a web page,
{\em share} is the fraction of those shoppers who buy Product $P$, and
{\em count} is how many of product $P$ they buy.
If we assume traffic is absolutely constant (so many per hour),
shoppers don't influence each other or share common influences,
share doesn't vary, and count is always one (1.00),
then the distribution of unit sales is Poisson
and our models are based on truth.

If average count is more than
one unit per shopper purchase (``market basket''),
then variance to mean will be a steady ratio.
We can call this a {\em first order} overdispersion.
Let's explore this one in a little more detail.
Suppose we have an average count $c > 1$.
Let's model count as shoppers
who buy product $P$ as having a fixed probability $r$
of buying one more.
Here are three logical check points for $c$ and $r$.
\begin{equation}
c=1 \iff r=0
\end{equation}
\begin{equation}
c=2 \iff r=1/2
\end{equation}
\begin{equation}
c \to \infty \ {\mathrm as} \ r \to 1
\end{equation}

This count model is a nice exponential model
with one or two obvious exceptions.
For example, there are wheel casters
for sofas and beds where the probability
of buying four given buying at least three
is a lot different than the probability
of buying five given buying at least four.
Batteries can be another example.
But, on the whole, the exponential count model is a good one,
and certainly a convenient one.
These probabilities add up to one ($1.0$).
\begin{align}
        Prob (X=1) &= (1-r) \nonumber \\
        Prob (X=2) &= r(1-r) \nonumber \\
        Prob (X=3) &= r^2(1-r)          \\
        Prob (X=4) &= r^3(1-r) \nonumber \\
        Prob (X=5) &= r^4(1-r)  \nonumber \\
                   &...          \nonumber
\end{align}
If $c=1$ (so $r=0$), then $Prob (X=0) = 1$ and all the rest are zero.
We're assuming, as a condition,
that this shopper bought at least one unit of this product,
so $Prob(X=0)=0$.

Remember this from math class in high school.
\begin{align}
1+r+r^2+r^3+r^4+... &= \sum_{i=0}^\infty r^i = \frac{1}{1-r} \\
1+2r+3r^2+4r^3+5r^4+... &= \sum_{i=0}^\infty i r^i = \frac{1}{(1-r)^2}
\end{align}

Let's look at the relationship between $r$
\begin{align}
c = E(X) = \sum_{X=1}^\infty
     &= \sum_{X=0}^\infty X (1-r) r^{X-1}                  \nonumber
     &= \\
%    &= (1-r) \sum_{X=0}^\infty \left( X-(X-1) \right)
%                               \sum_{i=X}^\infty r^i      \nonumber
%         \\
     &= (1-r) \sum_{X=0}^\infty \sum_{i=X}^\infty r^i
     &= \\
     &= (1-r) \sum_{X=0}^\infty r^X \sum_{i=0}^\infty r^i  \nonumber
     &= \\
     &= \frac{1}{1-r}                                      \nonumber
% \\ &= c \ \ \ \ \mathrm{(as\ expected)}                  \nonumber
\end{align}

%%%%%%%%%%%%%%%%%%%%%%%%%%%%%%%%%%%%%%%%%%%%%%%%%%%%%%%%%%%%%%%% <<<

Working out the conditional distributions in analysis
frustrated me, so I wrote a simulation.
Given a value $c$ it computes $r=(c-1)/c$,
samples one~million ($1000000$) Poisson distributions
of each mean from one-fifth $0.2$ through ten ($10.0$),
and simulates the count model
using random increment with probability $r$.
I get a consistent variance-to-mean ratio $\sigma^2/\mu$ of $2c-1$.
$c=1$ gives a straight Poisson with $\sigma^2/\mu = 1.0$,
$c=1.5$ gives  $\sigma^2/\mu = 2.0$,
$c=2.0$ gives  $\sigma^2/\mu = 3.0$, and
$c=5.0$ gives  $\sigma^2/\mu = 9.0$.

With a suggestion from a long-time friend,
I approached the math another way and it came out
the same as the simulation.
So let's look at that math.
For a given count $c$ we have $r=(c-1)/c$ and $c=1/(1-r)$.
We can decompose our unit-sales distribution
into the sum of distributions for each specific count $k$.
I'll note the following three sums.
\begin{align}
\sum_{k=1}^\infty     r^{k-1} &= \frac{1}  {1-r}               \\
\sum_{k=1}^\infty k   r^{k-1} &= \frac{1}  {(1-r)^2}           \\
\sum_{k=1}^\infty k^2 r^{k-1} &= \frac{1+r}{(1-r)^3}
\end{align}
For mean count $c$ we have a Poisson of mean $\lambda / c$
giving us an overall mean of $\lambda$ for unit sales.
\begin{align}
\mu    &= \frac{\lambda}{c} = \lambda \frac{1}{1-r}
       \\
\mu_k  &= \sigma_k^2 = \lambda \frac{1}{1-r} Prob(count=k)
                     = \lambda r^{k-1} (1-r)^2
          \\
\mathrm{TotalMean}  &= \sum_{k=1}^\infty k \lambda r^{k-1} (1-r)^2
                     = \lambda (1-r)^2
                       \left( \sum_{k=1}^\infty kr^{k-1} \right)
          \\
                    &= \lambda (1-r)^2 \frac{1}{(1-r)^2} = \lambda
                                                                     \\
\mathrm{TotalVar}   &= \sum_{k=1}^\infty k^2 \lambda r^{k-1} (1-r)^2
                     = \lambda (1-r)^2
                       \left( \sum_{k=1}^\infty k^2 r^{k-1} \right)
                                                                     \\
                     &= \lambda (1-r)^2 \frac{1+r}{(1-r)^3}
                      = \lambda \frac{r+1}{1-r}
                      = \lambda (2c-1)
          \\
\frac{\mathrm{TotalVar}}{\mathrm{TotalMean}} &= 2c-1
\end{align}

%%%%%%%%%%%%%%%%%%%%%%%%%%%%%%%%%%%%%%%%%%%%%%%%%%%%%%%%%%%%%%%% >>>

If traffic varies from day to day,
then we expect a factor increase in the standard deviation
of unit sales
{\em second order} overdispersion.
Let's go to Wikipedia to find the variance of a product.
If we have two distributions
$(\mu,\sigma^2)=(\mu_x,\sigma_x^2)$ and
$(\mu,\sigma^2)=(\mu_y,\sigma_y^2)$, then
\begin{align}
\mu_{xy} &= \mu_x \mu_y                                          \\
\sigma_{xy}^2 &= (\mu_x^2 + \sigma_x^2)(\mu_y^2 + \sigma_y^2)
               - \mu_x^2 \mu_y^2
\end{align}
If $X$ is a Poisson distribution,
then $\mu_x = \sigma_x^2 = \lambda$.
Let's suppose the multiplicative factor for traffic
$(\mu_y,\sigma_y^2)$ has mean $\mu=1$ and variance $s^2$.
Then we get the following for variance.
\begin{equation}
\sigma_{xy}^2 = (\lambda^2 + \lambda^2) (1+s^2) - \lambda^2 =
                 \lambda^2 (1+2s^2).
\end{equation}
Notice that this overdispersion maintains a constant
standard deviation ratio
so $\sigma / \mu$ depends only on $s$
while the variance ratio
 $\sigma^2 / \mu$ is linear in the mean value $\lambda$.

\section {OUT OF STOCK}
  \label        {STOCK}

In my old days of retail science (Khimetrics 2003-2008)
we estimated out-of-stock before doing demand modeling.
In other words,
we looked at the sales history for each product-store combination
and guessed which sequences of zero sales were truly zero sales
and which had the product unavailable
(not on the shelf or, more recently, not on the web page).
Given a Poisson distribution of mean $\lambda$,
the probability of a zero sample is $e^{-\lambda}$
and the probability of a particular sequence of $n$ time periods
being zero is $e^{-n \lambda}$.
This is the out-of-stock model from early days of retail science
and it's the out-of-stock model today.
The difference is the estimate of $\lambda$ in the early days
came from average unit sales by time period
with no deeper understanding
while, today at Clear Demand,
we use our demand model to estimate the same $\lambda$.
The consequence of doing out-of-stock estimation before modeling
is that the model depends heavily
on the zero-sales availability guesses.

In any case, we use a threshold probability,
default around $0.02$ (two~percent),
where we assume that any probability above that
is available, in stock,
and any probability below that is not available, out of stock.
As $e^{-4} = 0.02$ we get four units being that threshold.
If it sells ten units a day, then any zero is out of stock.
If it sells three units a day, then one no-sales day ($3.0$~units)
is presumed in stock
and two no-sales days in a row ($6.0$~units) are presumed out of
    stock.
If it sells 1.5 per day, then two day ($3.0$~units) are in stock
and three  ($4.5$~units) are out of stock.
None of this is terribly certain,
which is why I consider it a breakthrough
not to use this calculation in determining
estimates of model coefficients $\beta$ and $q$.

Overdispersion makes a mess of this calculation.
Consider a product with high average count or
stores with highly-varying traffic.
If the post-model out-of-stock calculation is wrong,
then okay, we can deal with that,
but I'm a lot more nervous when
a pre-model out-of-stock calculation is messing up
my elasticity-$\beta$ values.

\section {ELASTICITY INFERENCE}
  \label            {INFERENCE}

In the economics-textbook world
elasticity and $\beta$ values are easy to calculate
from a sales history universally rich in price and quantity changes
over a long, well-populated sales history.
The reality is that only a minority of products and stores
have enough sales history information to yield
a meaningful $\beta$ coefficient
and we're going to need some techniques and tricks
to generate the rest of the $\beta$ values.

The journey from maximum-log-likelihood {\tt ModelBeta} values
to universal model-group \fl{Beta} values is a long one
with many twists and turns.
There have been various alternative calculations of $\beta$
used in RSIM when the model is absent or insufficiently confident.
This is the one we're using by default at the moment.
We base the later calculations on {\em elasticity}
rather than $\beta$ to eliminate effects from different prices.

                \subsection {PriceBeta, ModelBeta, and MixBeta}
                                               \label {Mix}

We have a model that produces {\tt ModelBeta}
and we have said much about this model in this paper
and in other places.
This is well-established.

We also have a reverse engineered {\tt PriceBeta}.
For a given revenue-profit weight $\lambda$
we know that the optimum price $P$ is
\begin{equation}
P = \frac{1}{\beta} + \lambda C
\end{equation}
where $C$ is cost.
It doesn't take a deep understanding
of high-school algebra to solve for $\beta$.
\begin{equation}
\beta = \frac{1}{P + \lambda c}
\end{equation}
If we suspend disbelief and assume the best price was selected,
then this is the {\tt PriceBeta} $\beta$ we're going with.

We generate {\tt MixBeta} from {\tt PriceBeta} and {\tt ModelBeta}.
If there is no {\tt ModelBeta} for lack of sales-events information,
then the decision is easy, {\tt MixBeta} is just {\tt PriceBeta}.
If the statistical standard deviation of {\tt ModelBeta}
is too high, then we also use {\tt PriceBeta} as {\tt MixBeta}
In this section of the RSIM algorithm
too high is {\tt ModelBeta} times
the {\tt PriorWgtThreshHigh} parameter.
If the statistical standard deviation of {\tt ModelBeta}
is less than {\tt ModelBeta} times
the {\tt PriorWgtThreshLow} parameter,
then we use the {\tt ModelBeta} for {\tt Mixbeta}.
Anything in between is linearly interpolated between
{\tt PriceBeta} and {\tt ModelBeta} for {\tt MixBeta}.

%ob 1/(price-(lambda*cost))
%qb model sig/beta below PriorWgtThreshLow pb, above High ob,linear
    btwn

                \subsection {Elasticity Confidence}
                      \label           {Confidence}

Our maximum-log-likelihood model calculations described in
Sections~\ref {PRICE} through Sections~\ref {COMPRESSING}
are wonderful and cool and clever,
but they only work if there are price changes
in the price history.
Even those $Q_0$ and $\beta$ values are not 100\% reliable
because the supporting sales data may not be compelling.
For example, suppose there are two years priced at \$5.99
with one day with a lowered price of \$4.99.
Do we really believe the model $\beta$?
What confidence should we have in our coefficients?

The statisticians have a wonderful answer for us.
It turns out, by some streak of good statistical fortune,
the inverse of the log-likelihood second derivative (Hessian)
of a Poisson distribution whose mean is a linear function
is the covariance matrix of the coefficients.
Put in plainer English
if we take the inverse of the Hessian matrix,
called the Fisher information matrix (I think),
the one we used in the Newton-Raphson method
to find $Q_0$ and $\beta$,
the terms of that matrix are the covariance values
of the coefficients being solved.
That means the diagonal elements are the coefficient variances
and the square roots of those diagonal elements
are standard-deviation confidence intervals.
So we can take the
Hessian-inverse diagonal $\beta$-element's square root
to get a confidence interval for $\beta$.

The trouble is that this confidence interval
doesn't tell the whole story.
It answers the statistical question:
What is the nature of variability of $\beta$
given (1) the distribution of the samples is precisely Poisson
and (2) the mean of that Poisson distribution is exactly
$e^{q_0-\beta P}$ for all prices $P$
for {\em some} values of $q_0$ and $\beta$?
Given all the varying factors affecting shopper demand
based on price, it behooved us to come up with
more-realistic criteria for confidence in $\beta$
for a price family, area, customer type, and sales type.

For each of these, we define a minimum and maximum
to scale it from zero to one.

\BL
Statistical confidence.
It may not be the only measure of $\beta$ confidence,
but it should still be on this list.

\BL
Number of sales events.
This is the number of distinct dates added up for
all the products in the price family
all the locations in the area.

\BL
Number of distinct dates.
This is the number of different dates used
in all the sales events in the products and locations.
This is an artifact of how Marina calculated
the number of events and I don't see how it is
different from the number of sales events in a good way.

\BL
Number of price increases,
number of price decreases, and
combined (bi-directional) price-changes score.
We can count the price increases as one measure of activity,
the price decreases as another measure of activity, or
half of each added together as a combined measure.

\BL
Price deviation.
If we take the average deviation of price from the mean
we get a measure of the significance of price change.
Having one day a month of a lower price
and having two weeks each month at a lower price
will both give the same number of price changes
while the latter is far superior in measuring
response to price.

\BL
Age (recency).
How recent is the last price change?

\BL
Units deviation.
It's all moot if the unit sales don't vary or don't vary much.

Once we make all these measurements and combine them
into zero-to-one numbers, we can take those numbers
to given exponential powers and multiply them together.
There were tests using GridRunner and we ended up, for now,
with number of events, price deviation,
combined bi-directional price changes,
age, and units deviation each getting equal weight
and multiplying them all together to the first power.
Interestingly, the statistical confidence didn't make our list.

If this group price score is above
the {\tt PriorAvgThreshHigh} parameter,
then we use the {\tt ModelBeta} $\beta$
as calculated without concern.
If this score is below the {\tt PriorAvgThreshLow} parameter,
then we rely entirely on the following inferred-$\beta$ calculation.
Otherwise we linearly interpolate based on the group price score.

                \subsection {Price Family and Area Adjustments}
                     \label {FA_Adjustment}

We find an average elasticity for each model group,
that is $\beta$ times the average price.
We model each price family (F) and each area (A)
as having some incremental change
to the overall average elasticity.
For a tiny example suppose
price family $F_1$ has average elasticity $0.8$,
price family $F_2$ has average elasticity $0.9$, and
price family $F_3$ has average elasticity $1.0$.
Assuming we weight them equally
then our ``average average'' would be 0.9,
the $F_1$ difference would be $-0.1$,
the $F_2$ difference would be zero, and
the $F_3$ difference would be $+0.1$.
Given that we have a collection of not-necessarily-consistent
average elasticities for several model groups
(within a customer type and sales type)
with overlapping price families and areas,
and given that these elasticities are
likely to be overdetermined and
unlikely to be perfect consistent,
there is a resolution of these differences to be done.

One issue is the model elasticities are widely spread.
When there are only ten or twenty sales dates
for most products and stores
the {\tt ModelBeta} values can be low or high
and still be statistically confident.
That means elasticity differences can be larger or smaller
than what we would reasonably believe.
We deal with that by putting a arctangent S-curve
on the {\tt ModelBeta} elasticity with the
customer type and sales type average at the inflection point
so small differences from the average are pretty-much left alone
and larger differences are driven towards the average elasticity.
The {\tt Elasticity2Exponent}parameter
is the maximum difference in the exponent,
so setting {\tt Elasticity2Exponent} to 0.05
keeps the elasticities in this calculation
within five percent of the average
and the default of 0.333 keeps those elasticities
within forty percent of the average.

Sometimes it's easier looking at a few lines
of computer programming code.
% \begin{minipage}{\linewidth}
\begin{verbatim}

    x is model elasticity
    c is average elasticity (center of distribution)
    y is range-limited elasticity to return
    r is the Elasticity2Exponent range parameter
                /* Arctan limits are Elasticity2Exponent over two pi. */
    e = r/(PI/2.0);
                /* Arctan limits applied to log of elasticity. */
    a = log (x/c);
    b = atan (a/e)*e;
    y = exp (b)*c;
\end{verbatim}
% \end{minipage}
This function has the virtues of a unit derivative
when the elasticity $x$ is near the average elasticity $c$
and bounding the returned elasticity $y$ to $c$
times or divided by a factor of exp({\tt Elasticity2Exponent})
so it doesn't go crazy for small-sample-size models.

% ================================================================???

Let's start simple and work our way to the big picture.
We want to find differences
$X_{F_i}$ for each price family $i$.
(Pardon my casual notation here.)
The inconsistency residual would be
\begin{equation}
\label{eq_sum_sum}
r = \sum_G \sum_H \left( (X_{F_G}-X_{F_H}) - (E_G - E_H) \right) ^2
\end{equation}
where $G$ and $H$ are model groups
with price families $F_G$ and $F_H$,
$X_F$ is the difference variable for prime family $F$,
and $E_G$ is the elasticity of model group $G$.

The idea is to minimize the difference between the
$X$-difference and the $E$-difference.
\begin{equation}
r = \sum_G \sum_H X_{F_G}^2-2X_{F_G}X_{F_H}+X_{F_H}^2
                       + X_{F_G}(E_G - E_H)
                       + X_{F_H}(E_H - E_G)
                       + (E_G - E_H)^2
\end{equation}

We take the various partial derivatives
to generate the coefficients.
(I'm going to be lazy and not try to make my notation
 work for generating the matrix $A$ and the right hand side $B$.)
The $A$ matrix gets all the terms with products of two $X$ values,
coefficients are all unity ($1.0$)
and the $B$ vector gets all the terms with one $X$ value,
coefficients are all average-elasticity differences.
The matrix minimization
\begin{equation}
\mathrm{minimize} ~ ~ X^T A X + B X + C
\end{equation}
becomes the derivative equation
\begin{equation}
\label{eq_diff}
A X = -2 B
\end{equation}
and we know how to solve that (we hope).

Throw in the areas along with the price families
and weight the terms of the sums with factors depending
on their group price scores above and we get the code in RSIM
that sets all the $A$ and $B$ terms to zero
and increments their coefficients based on the
terms of the double-summation in
Equation~\ref {eq_sum_sum}.
The dimension of square-matrix $A$ and vector $B$
is the number of price families (F) plus the number of areas (A).

Since so many of the group price scores are zero,
the $A$ matrix is often very sparse
and an overconditioned problem becomes an underconditioned problem
with a singular matrix.
Requiring a minimum weight of some small number like $0.001$
mitigates that matrix near-singularity and gets more-consistent
solutions with small differences.

             \subsection {Integrating Price Family and Area
          Adjustments}
                  \label {Integrating}

If we apply these $X$ corrections
using the price families (F) and areas (A)
to the ``average average'' elasticity,
we get good prior estimates of elasticities
for each model group (G).
Before we apply these corrections,
we weight and limit them.
There are coefficients
{\tt AreaElasAvgWgt} and {\tt PFElasAvgWgt}
that can reduce the $X$ solution from Equation~\ref {eq_diff}.
There is another step that limits the deviation
from the average average elasticity to
{\tt AreaElasAvgRng} and {\tt PFElasAvgRng}
using arctangent S-curves.
These prior-based elasticities are
called {\em new elasticities} in RSIM.

For the model-group-level elasticity we deliver
we don't use the straight {\tt ModelBeta}
but the {\tt MixBeta} term
described in Section~\ref {Mix}.
We turn the group price score into a normal distribution $d$
as a sort-of confidence interval in beta $e$
and multiply by price to get an elasticity term $f$.
Note that the elasticity confidence bounds
are in terms of absolute elasticity rather than
fractions of anything.
We have higher confidence in terms with low absolute elasticity
confidence bounds rather than some fraction of the elasticity.
Based on the prices of each product in the price family
and each location in the area,
we compare the elasticity to the
{\tt ElastMin} and {\tt ElastMax} parameters
to keep in in range, typically between
$0.25$ and $4.0$.
This result is called {\tt Beta} or {\tt FcstBeta}.
(Later, when we have cannibalization and affinity terms
we'll product {\tt OptBeta} that takes those into account.)

Finally, there is a product-store model
with a separate maximum-log-likelihood model
for each product and store.
If the statistical standard deviation of that
product-store calculation of elasticity
based on the model group average price
is less than {\tt PriorWgtStoreLow},
then we use the product-store calculation of $\beta$.
If it's more than {\tt PriorWgtStoreHigh},
then we use the model-group {\tt Beta}
for the product and store.
In between we do linear interpolation based on the
elasticity standard deviation
and these two {\tt PriorWgtStore} parameters.

%          F_PARAM  (dpriorlow,  "PriorWgtStoreLow");
%          F_PARAM  (dpriorhigh, "PriorWgtStoreHigh");
%          F_PARAM  (epriorlow,  "PriorSmoothStoreLow");
%          F_PARAM  (epriorhigh, "PriorSmoothStoreHigh");

% Later on in RSIM's process flow
% we use cannibalization and affinity

\section {CANNIBALIZATION AND AFFINITY}
  \label {CANNIBALIZATION}

When one product's increase in sales
reduces sales of another,
we call it ``cannibalization.''
The example I use is putting Coca Cola on sale
and seeing Pepsi sales go down
as people presumably substitute the on-sale drink
for their usual choice.
(This example is a good story,
but I recall a study from my Khimetrics days
where the sale on Coke generated {\em more} Pepsi sales.
Maybe the promotional lift on all soft drinks
outweighed any cannibalization.)

The other interaction, in the opposite direction,
is ``affinity'' where sales of one product boosts
sales of another product.
My standard example is selling promoted hot dogs
increases sales of buns, catsup, and relish.
In that case the other products are not in the same
subcategory and, therefore, not in the same RSIM case
and RSIM would have no way of detecting that affinity.
There are, we presume, cases of affinity
within a single RSIM case.

So what do we measure?
More practically, what {\em can} we measure?
In some ideal world we would know cannibalization
and affinity for every ordered pair of products,
driver and dragged.
That's a lot to ask and what would that be based on?
With transaction-log (TLOG) data,
we might have a chance of seeing which pairs of products
tend to be bought together for affinity or,
conversely,
tend {\em not} to be bought together for cannibalization,
but RSIM doesn't have TLOG data and
there are no plans to add TLOG data to RSIM's input portfolio.

With just product-store-time-promotion sales data
and only within one subcategory,
we're going to have to set our sights lower.
Rather than trying to discern two different forces,
cannibalization and affinity,
we're going to consider a single vector
of positive attraction for affinity
and negative repulsion for cannibalization.
Further we are far short of enough data
with enough statistical precision
to evaluate this affinity-cannibalization coefficient
for every ordered pair of driver and dragged products.
At Khimetrics we retreated to a single coefficient $\psi$
for all cannibalization in our optimization case
that we called a ``micromarket.''
If all forecast sales went up 10\% and
cannibalization $\psi$ was 0.2,
then we would impose a 2\% downward adjustment
on all the forecasts in that micromarket.

In RSIM we're pursuing a middle path.
We calculate cannibalization-affinity on a single
product or price family as it relates to the entire
community of other products in the RSIM case,
usually a subcategory.
There are two ways to do this.
One is to ask how sales of other products are influencing
a particular product, good for forecasting,
not terribly useful for optimization.
Instead we ask how sales of one product
affect the sales of other products in the RSIM case.
We have price, promotion, seasonality, and availability
all affecting sales, but the question we ask in RSIM
for each product is this:
Is the {\em difference} between actual and model sales
of other products
correlated with the model or sales of this product?
If selling lots of Product $P$ causes other products
consistently to sell less than their models suggest
and lessing less of Product $P$ causes other products
to sell more than their models,
then we conclude that Product $P$ has a negative,
cannabalizing, effect on those other products.
We don't know {\em which} products are adversely affected,
only that there is a consistent negative-slope relationship
that suggests that selling more of Product $P$
takes sales away from other products
and reduces the total effective gain
of any price change or promotion.

That slope calculation is a simple linear regression
of a slope of $y$ as a function of $x$
where $y$ is the actual minus model of everything else
and $x$ is the average of actual and model of each product.
This value, negative for cannabalization
and positive for affinity,
appears in RSIM's output.

RSIM can optimize price based on the regular, plain-jane $\beta$
computed value based on maximum likeliood,
which we call {\tt FcstBeta}.
It can also optimize based on {\tt OptBeta}
based on {\tt FcstBeta} times
one minus cannabalization or one plus affinity
(they're the same calculation,
just different sign of the cannibalization-affinity-slope
regessson slope.
It's not an optimization panacea,
but it does directly react to the response
of other products to sales of Product $P$.

In general I lean towards using {\tt OptBeta} over {\tt FcstBeta}
for price optimization
even though there are some anomolies,
such as how did some point not the maximum-profit point
make more money on this product than the aximum-profit point?
I prefer using a measure of response to price $\beta$
that includes the effect of the price change
on other products as well.                                           

\section {RSIM DATA}
       \label {DATA}

The RSIM program does a lot of things
and requires a lot of different data to do those things.
Up until here in this paper we have been modeling
shoppers response to price and other aspects of nature
like seasonality, cannibalization, and affinity.

In keeping with my programming style,
I'll assign letters to my data primitives here
just as I do in my RSIM code.

I'm not following RSIM's data flow here,
I'm not following RSIM's history,
rather I'm following an expository evolutionary path
that shows how we look at the retail world in RSIM
one step at a time.
We start with the mathematical model in Section~\ref {Modeling}
(which was the sole purpose of this paper when I started writing it)
and proceed to the delivery
of price recommendations in Section~\ref {Pricing}.
From there we go to rules of various sorts in Section~\ref {Rules}
and finally to the opportunity curve in Section~\ref {Curve}.

                \subsection {Modeling Data}
                     \label {Modeling}

Already familier with the mathematical algorithms
from the first \ref {CANNIBALIZATION} sections,
let's look at the internal data arrays behind the math.

                            \subsubsection {P - Product}
                                    \label {P}

A product is something shoppers buy,
presumably different products are perceived as different
while things perceived as being the same are the same product
even if their ShopKeeping-Unit (SKU) or
Universal-Product-Code (UPC) numbers aren't the same.

                            \subsubsection {F - Price Family}
                                    \label {F}

A price family is a collection of products
that must have the same price in RSIM price recommendations.
Products in a price family are assumed to have the same
response to price ($\beta$),
but they can have different attributes and seasonality.
If the price family is not identified in a product input record,
then the product is considered to be in its own
{\em trivial} price family with just the one product.
We model by price families as explained in Section~\ref {G}.

                            \subsubsection {L - Location}
                                    \label {L}

A location is a store, a place of purchase, or a web page.
Actually, locations are the granularity of the sales-history
input data.
Since we have been doing business for clients,
locations have been stores,
but our original demonstration case for The Home Depot
had sales data in terms of markets, so locations were markets.

                            \subsubsection {A - Area}
                                    \label {A}

Areas are collections of locations,
usually as markets are collections of stores,
and we model areas as explained in Section~\ref {G}.

                            \subsubsection {T - Time Period}
                                    \label {T}

Time periods are either weeks are days for the sales history
and for the forecast interval.

                      \subsubsection {K - Customer Sales Type
          (CStype)}
                              \label {K}

We have customer types in Section~\ref {Cty}
and sales types in Section~\ref {Sty}.
(Before RSIM I wrote programs that used single letters
and, occasionally, single digits for my primitive data types.
When RSIM became large enough that I ran out of letters and numbers,
I decided it was okay to have multi-character primitives like
\fl{Cty} for customer type,
\fl{Sty} for sales type,
\fl{Pnl} for panel, {\em et cetera}.)
The combination of the two is a customers sales type (CStype)
and the letter \fl{K} comes from my recollection
of a similar construct at Khimetrics.

In Khimetrics lingo,
a CStype that could coexist with another sales type
based on customer type exclusivity or any other reason
was called a {\em parallel} CStype.

                                 \paragraph {Cty - Customer Type}
                                            \label {Cty}

Customer types are different populations of shoppers,
seniors, college students,
or the regular-sales, ``everybody'' population called {\em Regular.}
There are also self-selecting customer types
such as loyalty card membership.

                                 \paragraph {Sty - Sales Type}
                                            \label {Sty}

Sales types are the conditions of sale.
Normal, everyday, regular-price purchases are {\em Regular} sales
        type,
there are flyers, emails,
endcaps with the product at the end of the aisle,
signs in the store,
and signs on the fuel pump called ``pump toppers.''
Some sales types are only available for certain shoppers,
so they're tied to a customer type,
like a lower sale price for loyalty-card holders.

                            \subsubsection {G - Model Group}
                                    \label {G}

A model group is a combination of
a price family (\fl{F}),
an area (\fl{A}), and a CStype (\fl{K})
for modeling response to price $\beta$.
We abbreviate this combination as \fl{FAK}.

                            \subsubsection {E - Sales Event}
                                    \label {E}

A sales event is a report of product sold at a location
at a time with a CStype (\fl{PLTK}).
With parallel CStypes we can have two sales events
in the same product, location, and time (\fl{PLT}).
Even without parallel CStypes we can have two sales events
in the same product, location, and time (\fl{PTL})
from different parts of the time period,
especially if the time period is seven days
and a promotion doesn't last the entire week.

                            \subsubsection {Sgp - Season Group}
                                    \label {Sgp}

A season group is a collection of
products, locations, and CStypes (\fl{PLK})
that share common seasonality behavior.

                            \subsubsection {Sva - Season Variable}
                                    \label {Sva}

A season variable is one of the function-of-time variables
in the seasonality mix.
If $\theta$ is a time variable scaled
so one year adds $2\pi$ to $\theta$,
then our default Fourier season variables are
$\sin \theta$,
$\cos \theta$,
$\sin 2\theta$,
$\cos 2\theta$,
$\sin 3\theta$, and
$\cos 3\theta$
in the command script.
Seasonality is $e$ to a linear function
of these functions in time,
so these periodic trigonometric functions
have more up than down.
We have a linear-trend variable linear
which presents as an exponential function of time.

RSIM gets holidays from the {\tt Holiday} input files.

                            \subsubsection {X - Internal Data Matrix}
                                    \label {X}

When RSIM has to do some linear algebra
in maximum-likelihood modeling,
for putting good model groups together
to estimate bad ones,
and for reverse engineering attribute values,
we use these linear-algebra variables
to build matrixes and right hand sides
to solve these systems of equations.

                \subsection {Pricing Data}
                     \label {Pricing}

Pricing depends on a different aggregation
of products and stores than demand modeling.
We can use the same location clusters for pricing
that we used for demand modeling,
but we can use different pricing zones
that we call regions (\fl{Reg})
and we gather into decision prices (\fl{D}).

                            \subsubsection {Reg - Region}
                                    \label {Reg}

A region is a collection of products and locations
that are decreed to be assigned the same price.
Regions not only can split areas in the location space
they can split price families in the product dimension.

                            \subsubsection {D - Decision Price}
                                    \label {D}

A decision price is the aggregation entity
of products, locations, and CStype (\fl{PLK})
for RSIM to assign a price.
The scope of a decision price (\fl{D})
is price family, region, and CStype (\fl{FRegK}).

                \subsection {Rules Data}
                     \label {Rules}

    I have a lot to say about rules in Section~\ref {RULES}.
Since the primary message of Section~\ref {RULES} is how
rules work in RSIM,
it behooves us to spend some time now to establish
the data primitives we'll need to understand how RSIM works
on decision prices (\fl{D}).
Now we're only using hard rules with priority and rule rank
rather than soft rules with economic consequences of violation.

                            \subsubsection {B - Bounded Rule}
                                 \label {B}

Bounded rules involve only one decision price.
There are lower bound (possibly zero)
and upper bound (possibly max-price)
for each decision price (\fl{D})
along with a priority (higher is more important)
and rule rank (lower is more important if priorities are equal).

                            \subsubsection {Pnl - Panel}
                                 \label {Pnl}

Panels are collections of products and locations
with an assigned priority.
Rules and panels can be tied together
with parameters associated with the combination.

                            \subsubsection {Att - Attribute}
                                 \label {Att}

Attributes are text strings associated with
products, price families, locations, areas, and regions.
If a rule has a numeric field,
then that field can be ``indirectly'' associated
with an attribute.

                            \subsubsection {Cpi - Competitor}
                                 \label {Cpi}

Competitive rules are associated with specified competitors.

                            \subsubsection {Cpr - Competitive Price}
                                 \label {Cpr}

Competition is expressed as competitive prices
by competitor, product, store, date, and sales status
with the competitor and product being required.

                            \subsubsection {Sta - State}
                                 \label {Sta}

For tobacco state-minimum-price rules locations have states.

                            \subsubsection {R - Comparative Rule}
                                 \label {R}

Comparative rules involve more than one decision price.
Bounded and comparative rules are sufficiently different
that they're distinct data types in RSIM.
Rules with more than two decision prices
are called Initial Mark-Up (IMU)
and they use the same comparative rule data type (\fl{R}).

                            \subsubsection {C - Constraint}
                                 \label {C}

Comparative rules are expressed by a collection of
decision-price-to-decision-price (\fl{D}-to-\fl{D}) constraints.
These are linked-listed various ways
so RSIM can find them quickly.
Some of them go in only one direction to support one-way rules.

                            \subsubsection {J - Price Network}
                                 \label {J}

All the decision prices (\fl{D}) connected to each other
by one or more constraints form a price network.
Price networks (\fl{J}) are rule-and-elasticity optimized together.

%                            \subsubsection {x}
%                                \label {x}

                \subsection {Opportunity Curve Data}
                                 \label {Curve}

I designed RSIM to develop the revenue-profit opportunity curve
in two stages.
There is an optimization-points stage
with rules and ending-number price points
for varying revenue-profit weight ($\lambda$)
and then those optimization points are combined
to form display points for the opportunity curve output.
Ideally the optimization point calculated
for a given $\lambda$ value would be the
display point on the curve,
but heuristic searches around rules and ending-number price points
are uneven enough that sometimes the ``wrong-$\lambda$''
optimization points is used for a display point
and sometimes the $\lambda$ that looked good during optimization
isn't the best choice for the final display curve.

                            \subsection {Opt - Optimization Curve
          Point}
                                 \label {Opt}

For each price network (\fl{J})
optimization points are calculated
for several revenue-profit weights
in attempts to cover a range of $\lambda$ values
and to get a $\lambda$ value close to the
target profit-to-revenue margin
(usually the current-price margin).

                            \subsection {Dpt - Display Curve Point}
                                 \label {Dpt}

Once we have a collection of price-network-revenue-profit
(colloquially \fl{J}$\lambda$) optimization points

\section {RULES}
  \label {RULES}

Back in my old Khimetrics days (2003-2006)
price recommendation was about optimization
based on the price-elasticity model and
the ``Associated Product Rules'' (APR) were an afterthought.
When Clear Demand started working with our first client
we found out that rules were important enough
they came before elasticity-based optimization
in our work for them.
We started with competitive rules in Section~\ref {Competitive}
followed by step rules in Section~\ref {Step},
price rules in Section~\ref {Price},
and margin rules in Section~\ref {Margin}.
We can use attribute values for rule parameters
in Section~\ref {Indirect}.
In anticipation of getting business is soft-lines retail
with inventory management we developed
sell-through rules in Section~\ref {SellThru}.
The Home Depot had their own heuristics for
resolving competitive prices
which we have generalized in Section~\ref {CprResolution}
and prioritizing their rule alerts in Section~\ref {ConstraintBook}.

Early in our RSIM work we came up with
a structure for comparative rules in Section~\ref {Comparative}.
Comparative rules are guidelines that express themselves
as constraints (\fl{C}) between two decision prices (\fl{D})
% There are price rules in Section~\ref {EnumeratedR},
There are choice rules in Section~\ref {Choice},
measured rules (that we don't use) in Section~\ref {Measured},
and only-one-way ``equal-bang'' rules in Section~\ref {OneWay}.

Rules with more than two decision prices were once called
Initial Mark-Up (IMU) rules in Section~\ref {IMU}
and the name stuck.

Dealing with combinations and conflicts of bounded rules
is in Section~\ref {BResolution}
and the much-harder issue of dealing with
combinations and conflicts involving
comparative rules in Section~\ref {Resolution}
is Ever So Much More So.

Output files describing rule activity
are in Section~\ref {ConstraintBook}.

We'll start by going over the stronger-than-rules
price families and ending-number price points
and then we'll introduce the loci of coverage,
rule panels, in Section~\ref {Panels}.

                \subsection {Price Families and Regions}
                     \label       {Families}

While Sales input data can have different prices
for different products in a price family
or different locations in a region,
RSIM stores prices for each decision price.
There is no opportunity to have difference prices
within a price family or within a region.
It's not a matter of panel priority or rule rank.

There are several different prices for a decision price,
individual, current, elasticity, optimized, recommended, and export.
These prices may or may not be different from each other,
but they are the same for all product and locations
in the price family and region of the decision price.

There is an exceptional ``corner case''
where productds in price families can have
different prices from RSIM.
If products in a price family are in distinct regions,
then they can have distinct prices
because now they are in distinct decision prices.
The picture here the blue flavors of yogurt
like blueberry and grape come from one vendor
while the red flavors like cherry and strawberry
come from another vendor,
they have distinct costs and we can give them different prices.

                \subsection {Ending Numbers}
                     \label {Ending}

RSIM has price lattices of ending-number price points
so prices can end in nine like \$1.49, \$1.69, \$2.29, and \$3.89
or ninety-nine like \$1.99 and 3.99
or ninety-five like \$9.95 and \$11.95.
Products belong to price lattices
(and I don't remember what happens of two products
in one price family belong to different lattices).
These are {\em almost} as inviolate as
price families having the same price.

The only exception to price-lattice, ending-number pricing
is that a decision price is allowed to keep its current price
if the {\tt CurrentOK} parameter is turned {\tt ON}.
Current price is
the price derived by minimum mode calculations from
individual product-location or product-area prices.
If the decision price is currently \$5.55, for example,
then that price can be used for
elasticity price, optimized price, and recommended price.

The assignment of prices by RSIM
stays in this ending-number, price-point list
no matter what the other rules say.

                \subsection {Panels}
                     \label {Panels}

One of the more interesting aspects of retail rules
is defining the locus of coverage of a rule.
The scope of a rule can be the entire enterprise,
a collection of stores, specific products,
or some odd collection of products and locations.
A panel comes with a positive-integer priority
(higher number is higher priority)
and the rule-panel combinations comes
with a positive-integer rule rank
(lower number outranks higher number
with the same priority).

It only takes one product-location finger
of a panel touching a decision price to count in RSIM.
If a rule says a price has to be less than, say,
seven dollars \$7.00
and only one product and location in the
price family and region is in the panel for that rule,
then the rule still applies.

Panels allow all kinds of flexibility
beyond rule-assigned priority.
Picture a system that gives each competitor a priority,
so Competitor~A is always considered more important
than Competitor~B.
It may turn out Competitor~A is {\em usually}
more important than Competitor~B,
but for some products or some locations
Competitor~B is more important than Competitor~A.
Panels make it possible to have different priorities
for the same competitor for different products and locations,
so they can do this.

Panels facilitate multiple rules for the same competitor
with different parameters.
We can have a high priority panel keeping us
within 15\% of a competitor
but we would {\em like} to be
within 5\% of that competitor.
We can have multiple rules or panels
for the same competitor to create this kind of
tiered response to that competitor.
The same applies for multiple responses to margin
or step limitations.

Panels extend to comparative rules as well.
For two decision prices to be compared by a rule,
they both have to have a product-location combination
in the same rule and panel.

A rule can have many panels and
a panel can have many rules.
Panels are assigned to rules, products, and locations
in the RuleProductLocation input.

                                \subsubsection {Priority}
                                        \label {Priority}

Priority is associated with panels rather than rules.
It's a positive integer without any particular upper limit,
so clients often use numbers in the tens of thousands for
panel priorities.
The higher number is higher priority.
There is no advantage having a bigger gap
as 15 is just as much a higher number than 13 as 10000,
but there's no harm in using large numbers
for panel priorities.
(I sometimes wonder if the motivation behind clients
using five-digit integers for panel priorities isn't
the same as the group trying to prioritize their work
labeling all their projects ``URGENT.'')

                                \subsubsection {Rule Rank}
                                        \label      {Rank}

Rule rank is associated with rule and panel.
It is also a positive integer that runs the opposite direction
so the lower rule rank takes priority over a higher rule rank.
(``So priorities go up and rule ranks go down.
Don't you guys talk to each other?'')
So the most that rule-panel rule ranks can do, mostly,
is prioritize rules within the same panel
as two different panels, supposedly, never have the same priority.
There is an exception for competitive rules and
many of our clients assign the same priority to different panels
(with the same reckless abandon that some drivers have
when they drive five miles per hour over the speed limit).
In any event, the combination of panel priorities
and rule-panel rule ranks gives a clear ordering
of rule-panel combinations with ties allowed for competitive rules
and sometimes for other rules as well.

                \subsection {Bounded Rules}
                     \label {Bounded}

A bounded rule in RSIM is a rule defined for one decision price.
As the decision price has one price for
the entire price family, region, and CStype,
the bounds control
the entire price family, region, and CStype.

We have
competitive rules in Section~\ref {Competitive},
margin rules in Section~\ref {Margin},
step rules in Section~\ref {Step}, and
price rules in Section~\ref {Price}
with possible indirect prices in Section~\ref {Indirect}.
What makes life interesting here is the
minimum and maximum parameters in the RulePanel input
mean different things for all of these bounded-rule types.

Once we calculate a lower-bound price for
a rule and decision price,
we move that up to the next ending-number price point.
We call those prices {\tt Calc} and {\tt Bound}
in the output files in Section~\ref {ConstraintBook}.
We do the same thing with the calculated upper-bound price
we move that down to the next ending-number price point.

Sometimes the lowered upper bound ends up lower
than the raised lower bound, so there is no
ending-number-price-point-feasible price range.
Suppose a rule requires our price to be 70\%
of a competitor priced at \$8.99.
that puts us at \$6.29, presumably an ending-number price point,
and RSIM should ignore the \$0.003 left over.
If that same competitive rule requires our price to be 75\%,
then our rule-required price would be \$6.74.
If our nearest ending numbers are \$6.69 and \$6.79,
then our upper bound ends up \$6.69,
our lower bound ends up \$6.79,
and RSIM's answer to this conundrum
is to allow {\em both} ending numbers
by switching the bounds
so our upper bound is \$6.79
and our lower bound is \$6.69.

Many bounded rules can also have {\em target} prices
that are active when the rule has an upper or lower bound
and usually have the same format as the bounds
but with {\tt Target} instead of {\tt Maximum} or {\tt Minimum}
in the RulePanel input.
A target price is used if it is in a range of feasible prices.

                        \subsubsection {Competitive Rules}
                                \label {Competitive}

Competitive rules are what they seem,
price rules relative to a competitor.
These were the first rules we ran for a client,
the first rules in RSIM,
the rules we cut our teeth on at Clear Demand.
The principle is easy,
a competitor has a price and we want to maintain a price
not too much above and, maybe, not too much below their price.
We get those competitive prices by product and location
in Section~\ref {Cpr0} and
resolve those competitive prices for a decision price
in Section~\ref {CprResolution}.

Once we have a competitive price for a decision price
we can multiply
by {\tt Maximum} or {\tt Minimum} {\tt Factor} coefficients
and add {\tt Maximum} or {\tt Minimum} {\tt Amount} coefficients
to get our upper and lower bounds for the rule.
Targets are calculated from
{\tt Target} {\tt Factor} and {\tt Amount} fields.

We have different competitive rules for different competitors
but we can have different competitive rules with
the same competitor but different priorities, rule ranks,
or coefficients.

                            \paragraph {Competitive Prices (Cpr)}
                                                    \label {Cpr0}

We have a list, often a long list, of competitive prices,
usually gathered from competitors's web pages.
(I figure in the old days folks walked through
competitor stores with cameras or note pads writing down prices.
When I worked for Holiday Inn I asked a hotel owner
how he got the prices for the LaQuinta Hotel across the parking lot
and he said they called three times a day asking for rates
and he figured three calls he got every day were from them.)
These prices for a competitive item and store
are associated with one or more products and locations,
sometimes with a competitor sales type, regular or promotion,
sometimes with a date,
so there can be multiple competitor-price input records
with the same competitor, product, and location.

The {\tt IntelligentRuleFlag}
for a rule and panel
tells RSIM that the competitive rule
in this panel is only active in a specified elasticity range.
Picture Whole Foods selling a product for \$4.99
with a Kroger next door selling the same product for \$3.99
and yet the elasticity of that product at Whole Foods is low.
That suggests that the competition from Krogers to Whole Foods
isn't significant even though the two stores are close,
that people aren't price-comparing Whle Foods and Kroger
and, as a result, we ignore the rule under these conditions.


                                  \paragraph {CPR Calculation}
                                       \label {CprCalculation}

For each competitive rule, product, and location
we look for a competitive price.
The {\tt CpiPriceType} for the rule and panel
tells us whether we're looking for
regular price competition, promotion competition, or both
and the {\tt CpiHalfLifeDays} tells 
whether we're using just the most recent competitor price
or doing a weighted time average of the competitive-price history.
If {\tt CpiHalfLifeDays} is zero, then we use the most recent price.
Otherwise, we create a timeline of dated competitive prices
for the competitor, product, and location
and compute a weighted average from this history.


                                  \paragraph {CPR Resolution}
                                       \label {CprResolution}

Once each product-location that can have a competitive price has one
we put them all together for a single competitive price for the
decision price with all the product-location combinations (PLs)
in the price family and region.

If the PLs for the competitive rule
are in different panels with different priorities,
then we only use the PLs in the highest-priority panel
with competitive prices.

We started with the eight-twenty heuristic from The Home Depot.
The deal is that a price with 80\% of the PLs
is the competitive price,
then we use the lowest price that has at least 20\% of the PLs.
If there is no competitive price with at least 20\% of the PLs,
then the decision price has no competitive price at all.
Six locations with six different competitive prices will do that.
Those two percent numbers
(no hyphen, I don't mean 2\%,
I mean two numbers with percentages)
are global parameters
{\tt HighCompetitorThresholdPct}
and {\tt LowCompetitorThresholdPct}
so we can make a seventy-thirty rule or whatever we like.
Setting these two threshold-percent parameters the same
has RSIM using a minimum-mode for the competitive price.

                        \subsubsection {Margin Rules}
                                \label {Margin}

Margin rules are pretty simple.
We have a cost,
we usually have a minimum margin,
and we sometimes have a maximum margin.
Since margin is $(P-C)/P$ where $P$ is price and $C$ is cost
and since $C$ is the constant cost variable,
the price formula for margin factor $M$ is $C/(1-M)$.
Add a margin amount $A$ and the full formula is $P=(C/(1-M))+A$
where $P$ is price, $C$ is cost, $M$ is the margin fraction
in {\tt MinimumFactor} or {\tt MaximumFactor},
and $A$ is a constant amount
in {\tt MinimumAmount} or {\tt MaximumAmount}.
The same calculation with {\tt TargetFactor} and {\tt TargetAmount}
yields a target margin price.

The United States is actually one of few countries
that charge a price and then sales tax in addition.
(Unlike most countries,
taxes in the United States vary from place to place
so a nationally-advertised price would be awkward
with different tax rates in different places
while other countries typically have the same tax rate,
so that {\em could} be the reason for the difference.)
This price-included sales tax is often a value-added tax (VAT).
That tax is a fraction $\tau$ of the total selling price $P$
so the tax works out to $\tau P$ and it does {\em not} count
as part of the profit, so the profit is $P-\tau P-C$ and
the algebra gets us $P=(C/(1-\tau -M))+A$
with the extra tax term $\tau$.
This is especially convenient as the tax formula is the same
as the no-tax formula when $\tau$ is set to zero,
so this feature doesn't need any fork in the code,
just the extra $\tau$ term.

There is another wrinkle added to support a client,
United Dairy Farmers (UDF).
While they want profit and margin computed from the real cost,
they want the margin calculated from an attribute called
{\tt GrossCost} if that attribute is set.
We're still using the same cost in calculating profit
for RSIM output files,
just a different cost for the margin rule.
In addition they want the margin based on price minus
the buy down amount in the {\tt TotalBuydownAmt} attribute.
If we can use the {\tt GrossCost} attribute for cost $C$
and minus~{\tt TotalBuydownAmt} for the constant amount $A$,
then it all would work fine.

So how do we get these two terms into the margin equations?
First, subtracting {\tt TotalBuydownAmt} can be done using
indirect rules from Section~\ref {Indirect}.
Putting the attribute name in square brackets
{\tt [TotalBuydownAmt]}
in a calculated expression uses the attribute pertaining
to the product, price family, location, area, region,
{\em et cetera}.
Using gross cost instead of input-data cost
requires another twist because the input cost
is "hard coded" in the margin rule.
The factor and amount fields are already being used,
so we added a margin-rule feature where the RuleName
can be {\tt Cost=}$<$expression$>$ where the expression
can be {\tt [GrossCost]} to make the UDF version
of the margin rule work.

The UDF-style margin rule has {\tt Rule.RuleName} set to
{\tt Cost=[GrossCost]} and
{\tt RulePanel.}$<X>${\tt Amount}
set to {\tt -[TotalBuydownAmt]}
where $X$ is {\tt Minimum}, {\tt Maximum}, or {\tt Target}.
One test case had an attribute {\tt NegTotalBuydownAmt }
populated in {\tt ProductArea} input
with the same buy down amount as {\tt TotalBuydownAmt}
with a minus sign before it
with {\tt RulePanel.MinimumAMount}
and {\tt RulePanel.MaximumAMount}
set to {\tt [NegTotalBuydownAmt]}.

%
%
% 
%????????????????



                        \subsubsection {Step Rules}
                                \label {Step}

Step rules are a little more interesting.
Here the rule-panel {\tt Maximum} coefficients refer to the
outer limits, upper and lower rule bounds
while the {\tt Minimum} coefficients refer to the
smallest allowable price change from current price.
If the individual input prices for the product and locations
in the decision price are all the same,
then these do what one would expect.
Otherwise
the upper bound is computed from the lowest price
and the lower bound is computed from the upper price
so, if possible, no individual product and location
change more than the maximum step.
In case the individual price spread is too large
we never set the upper bound below the decision-price current price
and
we never set the lower bound above the decision-price current price.
Also, if individual input prices aren't all the same,
then there is no minimum step computed.

In case someone wants to allow more percent change up than down,
there are specific rule-panel parameters for
{\tt StepMaxFactor},
{\tt StepMaxAmount},
{\tt StepMinChangeFactor},
{\tt StepMinChangeAmount},
{\tt StepMinFactor}, and
{\tt StepMinAmount}.

Step rules have no target.

                        \subsubsection {Price Rules}
                                \label {Price}

The simplest rule of all is a rule that says
the price is between a lower bound and an upper bound,
for example
this product can't be less than \$2.99 or more than \$4.99.
It's easy enough,
just make the rule a price rule
and set the rule-panel {\tt MinimumAmount} for a lower bound
and set the rule-panel {\tt MaximumAmount} for an upper bound.
As presented here, that would require a separate panel
for each different price that product is allowed to have.
That may not be too bad
when there are two or three different prices,
but when there is a variety of products with lots of prices,
that would be tedious at best.
We can use indirect rules in Section~\ref {Indirect}
immediately below so we don't have to do all that.

{\tt TargetAmount} yields a target price.

                        \subsubsection {Indirect Rules}
                                \label {Indirect}

When a rule-panel coefficient should take on
different values and the use of different panels
to do that is awkward, unpleasant, or just plain ugly,
we can use RSIM attributes to assign values.
Instead of having a panel for each price or factor or whatever,
we can put the name of an attribute in brackets,
for example {\tt [CEILING]} and then assign the attribute
{\tt CEILING} in AttributeItemValue input
to whatever products, price families, locations,
areas, and even regions or panels as required.
This was a large community of different prices can be set
via a single price rule and panel.

                        \subsubsection {Sell-Through Rules}
                                \label {SellThru}

There is a full page http://radium/rsim/sellthru.htm
on sell-through rules, so I'll summarize here.
The term {\em sell through} refers to the fraction of inventory
sold in a time interval, so a sell through of 30\%
means three-tenths of the initial inventory is sold
between start and finish and seven-tenths remain.

For a single product and location
given start and end dates,
our model with seasonality answers the question
of what high price will sell the minimum
and what low price will sell the maximum fraction
of the given starting inventory.
We don't need availability since we know the inventory
and we know the product is in stock at the location.

For a collection of products and locations in a decision price
with different inventories and $Q$ rates of selling,
it's more complicated.
We could take the total units sold at each price
compared to the total inventory,
but some very slow movers would make achieving
the rule sell-through impossible.
(Math problem: We drive averaging thirty miles per hour
for half an hour, how fast do we have to drive to average
sixty miles per hour for a full hour?  Yes, like that one.)
We decided to weight the sales by the $Q_0$ factors so
slow movers don't count much and we can still achieve
our high and low objective at some prices.

                        \subsubsection {Bounded Rule Resolution}
                                \label             {BResolution}

If we have a collection of bounded rules
with lower and upper bounds
(possibly zero or infinity),
then how do we find a price or price range?

If we're lucky, then there is a range of feasible prices
at or above all the lower bounds and
at or below all the upper bounds.
When this happens there is room for price optimization.

When we're less lucky there is no feasible price,
no price
at or above all the lower bounds and
at or below all the upper bounds
and we're going to have to pick a price that violates
at least one rule.
We want to violate the lowest priority rules we can,
to violate higher priority rules by as little as we can,
and, in the event of a tie, to use the lowest price.

In this discussion priority takes into account
both panel priority and rule-panel rule rank.

We start with a blank interval from zero to infinity
and select the highest priority bounded rule.
In that selection process,
if the next bounded rules are of the same priority,
then we choose the rule with the lower upper bound.
(It doesn't matter which one has the higher lower bound.)
If the two intervals overlap, then we take
the higher lower bound and the lower upper bound.
If the new rule interval is completely below
the existing interval, then we reduce our interval
to just the existing interval's lowest point.
If the new rule interval is completely above
the existing inverval, then we increase our interval
to just the existing interval's highest point.
We repeat the process until
the upper and lower bounds are the same price
or we run out of rules.

Since all the intervals already take ending-number price points
into account in selecting their endpoints,
the final result of this resolution
will be a valid ending-number price point.
% ending numbers

                        \subsubsection {Target Price}
                                \label {Target}

% !!! put targets in for all rules that have them (not Step)

The {\em compliant} price is the price that follows
all the rules the best we can do.
If the interval for a decision price is a single price,
then it's easy to decide what the compliant price should be.
Otherwise we would pick the price in the interval
closest to the current price.

In the event that there are rules with targets,
we choose the highest-priority target within the range,
even if the range bounds have nothing to do with
the rule whose target we use.

                \subsection {Comparative Rules}
                     \label {Comparative}

The next step up the rules-complexity ladder
is a comparative rule,
a linear or affine relationship between two decision prices,
$A \le FB+G$ where $A$ and $B$ are prices of two decision prices,
$F$ is a positive-number multiplier factor,
and $G$ is a constant price amount.
These individual-decision-price-pair relationships
are called {\em constraints} in Section~\ref {Constraints}.

These come from various places.
There are constraint-described rules in Section~\ref {ConstraintInput},
choice rules in Section~\ref {Choice}, and
measured rules in Section~\ref {Measured}.
While measured rules are cool,
especially with the units-of-measure (UOM) support
so we can compare ounces-cups-pints-liters-gallons,
grams-ounces-pounds,kilograms, and units-dozens,
to my knowledge nobody has ever actually
used a measured rule in RSIM.

There is an interesting concept we have in RSIM
of a one-way rule where Price~A affects Price~B
but Price~B has no effect on Price~A.
The example was a big-city price affecting small town prices
of the same product
but a competitor in the small town lowering its prices
isn't going to propagate back to the big-city price.
We discuss these one-way rules in Section~\ref {OneWay}.
They go by the name ``equal-bang'' rules 

                                \subsubsection {Individual Constraints}
                                        \label            {Constraints}

A comparative rule is an umbrella.
Pairs of decision prices within that comparative rule
sharing the same panel can become constraints where
$A \le fB+g$ where $A$ and $B$ are prices of the decision prices,
$f$ is a positive multiplier factor,
and $g$ is a constant price amount.

                                \subsubsection {Constraint Input}
                                        \label {ConstraintInput}

Constraint input describes comparative constraints directly.
We specify a product, price family, area, or region in a CStype
compared to another product, price family,
location, area, or region in CStype
with a multiple and additive amount between them.
$A \le fB+g$ where $A$ and $B$ are decision prices,
$f$ is a multiplier factor, and $g$ is a constant price amount.
Whatever product, price family, location, area, and region
is not specified
is assumed to be all of them
and whatever customer type or sales type is not specified
is assumed to be regular sales.
(Writing the RSIM code to interpret these constraint input records
with all these combinations was not easy.)

For example,
if the constraint input compares Products $P1$ and $P2$,
then any two regions touching $P1$ and $P2$ in the same location
in regular CStype
are compared using the {\tt Multiple} and {\tt Add} values
$f$ and $g$.
A missing {\tt Multiple} is assumed to be one ($f=1.0$)
and a missing {\tt Add} assumed to be zero ($g=0.0$).

For another example,
if the constraint input compares Areas $A1$ and $A2$,
then any two regions in areas $A1$ and $A2$ that share a product
in regular CStype are compared
using the {\tt Multiple} and {\tt Add} values $f$ and $g$.

To facilitate a more-natural presentation
of constraint specification,
this input allows two groups of constraints stated over three 
products, price families, areas, regions, or CStypes.
The pattern is $f_1 P_1 + g_1 \le f_2 P_2 + g_2 \le f_3 P_3 + g_3$
where $P_1$, $P_2$, and $P_3$ are prices specified somehow
and $f_1$, $g_1$, $f_2$, $g_2$, $f_3$, and $g_3$ are coefficients.
The example is that the brand name Lucky Charms $P_2$
has to be between 25\% and 40\% more expensive
than the off-brand Marshmellow Maties $P_1$, also $P_3$,
so we have $1.25*P_1 \le P_2 \le 1.4*P_3$.
We could write this in two records,
$1.25*P_1 \le P_2$ and $P_2 \le 1.4*P_3$,
but the continuity is better represented in one longer record.
(I'm torn sometimes by the schizophrenetic distinction
between a data file being a machine-to-machine communication of data
and the same data file being intended to be read by humans
and even, occasionally, edited and written by humans.)

While we humans see structure when the first and third column groups
are the same
products, price families, locations, areas, regions, and CStypes,
RSIM doesn't care.
It processes the first and second groups of columns
and then does a shift-left
to process the second and third groups of columns.

In my effort to make this input totally flexible
for all kinds of constraint inputs,
I do require that the inputs be parallel,
whatever product, price family, location, area, region, and CStype
is or are specified in the first group of columns
({\em i.e.} {\tt ProductId1})
are specified in the second and third groups of columns
({\em i.e.} {\tt ProductId2} and {\tt ProductId3}).
It matters whether the
products, price families, locations, areas, regions, or CStypes
are the same or different,
but not whether one is there and another is missing.

I tried to handle all the cases of
presence or absence of
products, price families, locations, areas, regions, and CStypes.
Missing
product, price family, location, area, or region
is presumed to refer to all
products, price families, locations, areas, or regions
while missing CStype is presumed to be regular sales.
The idea is to find every combination of decision prices (D)
(price family, region, and CStype)
relating the two groups of columns
and creating individual D-to-D constraints
with the scale and slope calculated
from the $f$ and $g$ coefficients in the input file.

(To be honest, I'm not sure if RSIM actually does this,
but here are the plans for building constraints from rules.)

If just products are specified, $P1$ and $P2$,
then every location $L$ is checked for decision prices $D1$ and $D2$
(by price family, region, and CStype)
where $D1$ has product $P1$ at location $L$
and $D2$ has product $P2$ at location $L$
with $D1$ and $D2$ in the same CStype.
Those combinations $D1$ and $D2$ get individual constraints
associated with the rule.
This isn't entirely stupid if $P1$ and $P2$
are in the same price family
as a region boundary could put them
in two decision prices in some stores.

If just price families are specified, $F1$ and $F2$,
then we look at each location $L$
for products $P1$ in $F1$ in $D1$
and $P2$ in $F2$ in $D2$.
at each location to create constraints.

If just locations are specified, $L1$ and $L2$,
then every product $P$ is checked
for decision prices $D1$ and $D2$
(by price family, region, and CStype)
where $D1$ has location $L1$ with product $P$
and $D2$ has location $L2$ with product $P$
with $D1$ and $D2$ in the same CStype.
Those combinations $D1$ and $D2$ get individual constraints
with $D1$ and $D2$ associated with the rule.

Similarly they can specify areas $A1$ and $A2$
and we search for each product $P$
over all locations $L1$ in $A1$ with $P$ and $L1$ in $D1$
and $L2$ in $A2$ with $P$ and $L2$ in $D2$
to create constraints with $D1$ and $D2$.

We can specify both product and location dimensions,
preferably either the same product or the same location.
I suppose one could compare product $P1$ in location $L1$
to product $P2$ in location $L2$,
but methinks it makes more sense
to compare two products $P1$ and $P2$
or price families $F1$ and $F2$
at one location $L$ or area $A$ or even region $Reg$
or to compare two locations $L1$ and $L2$
or two areas $A1$ and $A2$ or even two regions $Reg1$ and $Reg2$
with the same product $P$ or price family $F$.

Specifying two CStypes $K1$ and $K2$ compares them
across every product and location
or those limited by those specified.
I'm presuming
the products, price families, locations, areas, or regions
would be the same.
I don't know if I should write RSIM code
for the case where one product or store on regular price
is compared to another product or store on sale.
Would anybody really {\em do} that?

                                \subsubsection {Choice Rules}
                                        \label {Choice}

We started our rules in RSIM with choice rules.
Consider four products,
$A$ small-brand-name,
$B$ large-brand-name,
$C$ small-house-brand, andf
$D$ large-house-brand.
\begin{minipage}{\linewidth}
\begin{verbatim}

                |
                |        A        B
        brand   |
                |
                |        C        D
                |
                +-------------------------------- size

\end{verbatim}
\end{minipage}
Consider
$A$ 12-ounce Lucky Charms from General Mills,
$B$ 18-ounce Lucky Charms from General Mills,
$C$ 12-ounce Marshmellow Maties (house brand), and
$D$ 18-ounce Marshmellow Maties (house brand).
Suppose the prices are
$A=\$3.89$,
$B=\$4.89$,
$C=\$2.59$, and
$D=\$3.69$.
The premiums for size are $B/A=1.26$ and $D/C=1.42$ and
the premiums for brand are $A/C=1.50$ and $B/D=1.33$.
Both size and brand are out of whack.
We could keep the same general price levels
and narrow those gaps by setting prices
$A=\$3.79$ (down ten cents),
$B=\$4.99$ (up ten cents),
$C=\$2.69$ (up ten cents), and
$D=\$3.59$ (down ten cents).
Now the premiums for size are $B/A=1.32$ and $D/C=1.33$ and
the premiums for brand are $A/C=1.41$ and $B/D=1.39$.

So all we have to do is fill in the size ratio of 1.33
and the brand ratio of 1.4 and it all works.
The problem is that coming up with those numbers
for the choice rules is about bringing price ratios
towards the average.
With only one rule the average might not be too hard
to estimate by eye,
but when two or more rules are involved,
RSIM does a smoothing calculation so each choice has a value
with the prices are close to their old values
except the choice ratios are more uniform.

% ???????????????? used to try to make a linear regression
% ???????????????? and matching every pair of decision prices
% ???????????????? looking for pairs that are otherwise the same
% ???????????????? and are in the same panel

                                \subsubsection {Measured Rules}
                                        \label {Measured}

Nowhere in the previous-section discussion
was there any cognizance of the size ratio
12~ounces to 18~ounces, one to 1.5.
The concept here is to force appropriate economy
for larger sizes of the same stuff.
We do that with an exponent assigned by rule and network.
The network here is a product network,
the sort of network we haven't used in RSIM
for a long time, if ever.
If we decide to do measured rules for real,
then we'll have to do something better
with this network dimension.

The idea is that we want more-quantitative control
than a choice rule gives us.
Let's say we want a factor of three in size
to give a factor of two in price.
If the 10-ounce can is four dollars \$4.00,
then the 30-ounce can should be eight dollars \$8.00.
Following the same curve,
15~ounces should be \$5.17,
20~ounces should be \$6.19, and
25~ounces should be \$7.13.
These points all follow a curve
\begin{equation}
\label{eq_exp_factor}
R_P = (R_S)^f
\end{equation}
where the price ratio $R_P$
is the size ratio $R_S$ to some given exponent factor $f$.
In this example the factor is 0.631 because $3^{0.631}=2$.

A simpler example would be a square root relationship
where $f$ is one~half (0.5).
If the size goes up by a factor of four, then the price doubles.
If the size goes up by a factor of nine, then the price triples.
If the size goes up by a factor of $x$,
then the price goes up by $\sqrt{x}$.

The idea here is to control the price increase as
size or weight or volume increases.
If the factor $f$ is one (1.0),
then there is no discount for a bigger size.
At the other end of the spectrum,
if the factor $f$ is zero (0.0),
then all sizes are the same price.
Normally we expect $f$ not to be too much less than one
as we can't afford to give the extra size away,
but one example stands out.
Cable ties were five~dollars \$5.00 for ten,
six~dollars \$6.00 for one~hundred,
and just over seven dollars \$7.20 for one~thousand.
The factor there was a paltry 0.079.
I guess somebody figured it was reasonable
to charge about five~dollars just to get a small number
and very few shoppers would pay much for many more.

As with the choice rules,
just asking the user to supply the factor $f$
is substituting one unfathomable value for another,
but we can run a statistical fit
to see what factor $f$ best fits their existing prices.

There is the issue of units of measure.
I remember being in a store where two brands of toothpaste
were "unit priced,"
one in cost per ounce and the other in cost per gram,
technically correct but not very useful.
(If I could easily multiply or divide by 28.3 in my head,
then I would easily be able to do the ratio calculation
for unit-price comparison and not need those unit prices at all.)
Similarly, we have input-date sizes in different input-data units,
some in pounds and ounces, some in quarts and gallons,
and even some in packs and ten-pack cartons.
With its units-of-measure entries RSIM can make the conversions
so the burden on those creating input data is lessened.

There is are differences in perception
between sixteen ounces and one pound
or between either of them and 453~grams,
but, for now, we'll just consider comparing the magnitudes
of heights, weights, or volumes to calculate
price ratios for rules.

                                \subsubsection {One-Way (Equal-Bang)}
                                        \label {OneWay}

%????????????????

                \subsection {Initial Markup (IMU)}
                     \label                 {IMU}

%IMU is R

                \subsection {Resolution}
                     \label {Resolution}

% subsection Optimization Search
% PARALLEL HIGHEST PRIORITY

                \subsection {AlertBook and ConstraintBook}
                                   \label {ConstraintBook}

%                \subsection {}
%                     \label {}

%                \subsection {}
%                     \label {}

\section {OPTIMIZATION}
  \label {OPTIMIZATION}

\section {OPPORTUNITY CURVE}
  \label             {CURVE}

                \subsection {Profit Weight (Lambda)}
                     \label                {Lambda}

                \subsection {Sequence of Curve Points}
                     \label {Sequence}

    Here's the deal.
We can select a profit weight ($\lambda$)
and apply it to price network.
A price network is a collection of decision prices
connected by comparative rules.
The problem is we don't know where the new curve point
for that $\lambda$ is going to end up until we
do the calculations.
The curve point is a piecewise-constant function
that jumps from one revenue-profit value to another
at certain $\lambda$ values.

    We're trying to find a $\lambda$-optimal price assignment
that has the same margin
(profit divided by revenue, proxy for price image)
as the currrent prices.

    We're also trying to get a curve of $\lambda$-optimal
points in revenue-profit space
so a more-energetic client can select
$\lambda$-optimal prices with a different objective coefficient.
A good curve has lots of points without
too much space between them,
what we would think of as a "pretty" curve.
It would be nice if that curve goes
from maximum revenue ($\lambda=0$)
to maximum proft ($\lambda=1$).
If the current-price margin is not in the range
$\lambda=0$ to $\lambda=1$,
then a "pretty" curve would extend to cover
that current-price margin.

    We can hope for a nice, "pretty" curve
with reasonably evenly-spaced points with only small gaps
spreading from $\lambda=0$ to $\lambda=1$,
maximum revenue to maximum profit,
and beyond that range if it is needed
to reach current-price margin.
It may not be possible to achieve a smooth curve
because the revenue-profit $\lambda$-optimum
may jump further than we would like
with a tiny, infinitesmal change in $\lambda$.
Also, especially if there is significant revenue-profit movement
due to rule compliance compared to more-lax current prices,
the $\lambda$-optimal curve may go nowhere near
current-price margin.

    Let's define a gap between two adjacent, nearest points
based on the revenue difference $\Delta_{rev}$,
the profit difference $\Delta_{pro}$, and
the maximum of the two revenues $Rev_{max}$.
\begin{equation}
\label{eq_exp_factor}
gap = \frac {( \Delta_{rev}^2 + \Delta_{rev}^2 )}{Rev_{max}}
\end{equation}


    Each point comes from a $\lambda$ chosen by our algorithm
based on curve points and their $\lambda$ values so far.

1. Let's do all the $\lambda$ values offered in an
{\tt OpCurveWeight} file if there is one.

2. If zero isn't on the list so far,
then the next choice is $\lambda=0$.

3. If one isn't on the list so far,
then the next choice is $\lambda=1$.

                \subsection {Selection of Curve Point}
                     \label {Selection}


\section {CONCLUSION}
  \label {CONCLUSION}

\end{document}
